\documentclass[10pt,twocolumn,letterpaper]{article}

\usepackage[review]{cvpr}
\usepackage{amssymb}
\captionsetup{font=footnotesize}

% ============================================================================
% numbers.tex -- every number that appears in the paper, with its source.
% RULE: no numeral enters main.tex directly. It is defined here, and the
% comment above it names the experiment directory that produced it. A number
% with no experiment behind it is written as \TODOnum so it cannot be mistaken
% for a measurement.
% ============================================================================

\newcommand{\TODOnum}{\textcolor{red}{\textbf{[unmeasured]}}}

% --- E00: DSEC exposure survey (experiments/e00_exposure_survey) -------------
% 18 train sequences, 21142 frames, measured from the published
% <seq>_image_exposure_timestamps_left.txt files.
\newcommand{\dsecExpMin}{118}              % us, zurich_city_06_a
\newcommand{\dsecExpMax}{14\,996}          % us, the auto-exposure ceiling
\newcommand{\dsecExpRatio}{127}            % max/min, rounded
\newcommand{\dsecFramePeriod}{50}          % ms, every sequence
\newcommand{\dsecPinnedFrac}{32.5}         % percent of frames at the ceiling
\newcommand{\dsecPinnedCount}{6866}
\newcommand{\dsecFrameCount}{21\,142}
\newcommand{\dsecCeilingSeqs}{six}         % sequences pinned at the ceiling
\newcommand{\dsecOtherCeilingSeqs}{five}   % the ceiling sequences not measured here
\newcommand{\dsecWithinSeqMin}{337}        % us, interlaken_00_d
\newcommand{\dsecWithinSeqMax}{4207}       % us, interlaken_00_d
\newcommand{\dsecWithinSeqRatio}{12.5}
\newcommand{\dsecCotriggerFrames}{7080}    % frames checked, L/R start identical
\newcommand{\dsecCotriggerPct}{100}
\newcommand{\dsecMidDiffMax}{190}          % us, max |mid_L - mid_R|
\newcommand{\dsecWidthDiffMax}{380}        % us, max |w_L - w_R|
\newcommand{\dsecTsResidual}{1}            % us, |image_ts - avg(mid_L,mid_R)|

% --- E01: events inside one exposure (experiments/e01_intra_exposure_change) -
\newcommand{\dayExpUs}{1478}               % interlaken_00_c, median
\newcommand{\dayEvents}{20\,616}
\newcommand{\dayPixFrac}{6.0}              % percent
\newcommand{\dayCrossings}{1.10}
\newcommand{\dayFrames}{536}
\newcommand{\nightExpUs}{14\,996}          % zurich_city_09_a
\newcommand{\nightEvents}{310\,320}
\newcommand{\nightPixFrac}{38.2}           % percent
\newcommand{\nightCrossings}{2.60}
\newcommand{\nightFrames}{452}
\newcommand{\nightDayEventRatio}{15.1}

% --- E02: RVT window alignment (experiments/e02_rvt_window_alignment) --------
% Read from the released source, not from the paper's description of it.
\newcommand{\rvtWindow}{50}                % ms, stacked_histogram_dt=50
\newcommand{\rvtBins}{ten}                 % bins per window
\newcommand{\rvtBinWidth}{5}               % ms per bin
\newcommand{\rvtUniformCentroid}{-25}      % ms, uniform-weight reference point,
                                           % superseded as a claim by E17.

% --- E04: RVT loads under torch 2.7.1 (experiments/e04_rvt_port_risk) --------
\newcommand{\rvtParams}{4.41}              % M, rvt-t
\newcommand{\rvtMissingKeys}{0}
\newcommand{\rvtUnexpectedKeys}{0}

% --- E06: DSEC-Det label forensics (experiments/e06_dsec_det_label_forensics)-
\newcommand{\ddBoxes}{390\,118}
\newcommand{\ddTracks}{6693}
\newcommand{\ddSeqs}{60}
\newcommand{\ddSpacingMedian}{50\,001}     % us
\newcommand{\ddSubFrameFrac}{0.0000}       % fraction of spacings below 40 ms
\newcommand{\ddSubFrameFracPct}{0.00}      % the same, as a percentage
\newcommand{\ddInterpResidual}{0.707}      % px, median second-difference
\newcommand{\ddInterpResidualPninety}{1.819}
\newcommand{\ddInterpBelowThresh}{7.1}     % percent below 0.01 px
\newcommand{\ddTriples}{361\,628}

% --- E07: LET-3D-AP citation graph (experiments/e07_let3dap_citation_graph) --
\newcommand{\letCitations}{40}             % returned by the API, coverage unknown
\newcommand{\letTemporalHits}{6}
\newcommand{\letAxisHits}{0}               % none apply the decomposition to time

% --- Removed 2026-09-05: quantities awaiting detection runs this paper does not
% contain. No detection score is computed anywhere in the paper or the
% supplement, so the tolerant-metric and re-anchoring macros below have no text
% that could use them and are left undefined rather than rendering a placeholder.
% \newcommand{\tauHatCI}{\TODOnum}
% \newcommand{\sigmaTau}{\TODOnum}
% \newcommand{\panelCSlope}{\TODOnum}
% \newcommand{\panelCIntercept}{\TODOnum}
% \newcommand{\apSyncDelta}{\TODOnum}
% \newcommand{\apBiasDelta}{\TODOnum}

%% ===== MEASURED 2026-09-02, CORRECTED 2026-09-03 ============================
%% E09+E11 per-object event-centroid dispersion inside one exposure. (E12 was
%% retracted on 2026-09-05, every value in this block is E11's.)
%% No detector, no checkpoint, no velocity denominator, no assumed center noise.
%% The values below are E11's exact-window re-run (result_exact.json), which
%% restricts events to [a,b] in microseconds. The first run sliced to
%% whole-millisecond boundaries and is superseded, it produced 208.8 us.
\newcommand{\sigmaEvtNight}{206.6}   %% us within-frame sd, night, all pixels
\newcommand{\sigmaEvtNullNight}{93.0}   %% us analytic uniform null
\newcommand{\sigmaEvtExcessNight}{183.6}   %% us excess over the null
\newcommand{\sigmaEvtFramesNight}{1134}   %% frames measured
\newcommand{\sigmaEvtFracNight}{0.9012}   %% fraction of frames sd>null
%% E11 flicker controls, exact window (result_exact.json)
\newcommand{\excessClean}{180.6}   %% us phase-locked pixels excluded
\newcommand{\excessRandCtl}{182.9}   %% us matched random exclusion
\newcommand{\lockedSensorPct}{4.67}   %% percent OF THE SENSOR, not of live pixels
%% E12 RETRACTED 2026-09-03 and removed from the paper 2026-09-04. An integer
%% window cancels a periodic source's zeroth moment, not its first: the centroid
%% contribution is -(A/lambda_0) sin(phi)/omega. The widths below were measured
%% correctly, the cancellation argument they were reported under was false, so
%% the macros are left undefined and no width appears in the paper.
%% \newcommand{\excessHalfPeriod}{212.6}   %% us w=5000 (0.5 periods)
%% \newcommand{\excessOnePeriod}{195.2}   %% us w=10000 (1.0 periods)
%% \newcommand{\excessFullExp}{183.6}   %% us w=14996 (1.4996 periods)
%% \newcommand{\excessTwoPeriod}{246.8}   %% us w=20000 (2.0 periods)
%% E08 label noise whiteness
\newcommand{\labelAcfMeasured}{-0.3566}   %% D3 acf lag1 measured on real labels
\newcommand{\labelAcfWhite}{-0.740}   %% same code on white noise
\newcommand{\labelAcfTheory}{-0.75}   %% analytic white-noise value
\newcommand{\ntracksNoise}{5250}
\newcommand{\nthirdDiff}{629\,442}
%% E10 mains flicker
\newcommand{\flickerLineHz}{100}   %% Hz, strongest night line
\newcommand{\flickerZmedian}{2.935}   %% night 100Hz Rayleigh Z, null 0.693
\newcommand{\flickerZnull}{0.693}
\newcommand{\flickerFracLocked}{0.2407}   %% fraction of LIVE pixels p<1e-3
\newcommand{\flickerCtrlNightZ}{0.581}   %% night 137Hz off-frequency control
\newcommand{\flickerCtrlDayZ}{0.719}   %% day 100Hz control

%% --- E00, shape of the exposure distribution (21142 train frames) ----------
\newcommand{\dsecExpGapLo}{4207}           % us, largest non-ceiling exposure
\newcommand{\dsecExpGapHi}{14\,996}        % us, the ceiling
\newcommand{\dsecExpGapCount}{zero}        % frames strictly between the two
\newcommand{\pctSixtySix}{66}              % the percentile index itself
\newcommand{\pctSixtyEight}{68}
\newcommand{\pctBandLo}{5}                 % Fig. 1's whisker band
\newcommand{\pctBandHi}{95}
\newcommand{\dsecExpPsixtysix}{2999}       % us, percentile 66
\newcommand{\dsecExpPsixtyeight}{14\,996}  % us, percentile 68

%% --- E06, the per-box time field of the released labels --------------------
\newcommand{\ddDistinctT}{70\,379}         % distinct values of t, summed within sequence
\newcommand{\ddBoxesPerT}{5.54}            % boxes per distinct value

%% --- E09, the measurement and its strata -----------------------------------
\newcommand{\sigmaEvtObjNight}{6}          % median qualifying objects per frame
\newcommand{\sigmaEvtSdPninetyNight}{344.5}% us, p90 of the within-frame sd (exact window)
\newcommand{\sigmaEvtPtpNight}{563.9}      % us, median within-frame peak to peak
\newcommand{\sigmaEvtDay}{18.0}            % us, interlaken_00_c within-frame sd
\newcommand{\sigmaEvtNullDay}{21.6}        % us, its analytic null
\newcommand{\sigmaEvtFramesDay}{20}        % usable frames, day
\newcommand{\sigmaEvtFracDay}{0.35}        % fraction of day frames sd>null
\newcommand{\excessDay}{0}                 % us, the day excess: sd is below the null
\newcommand{\sigmaEvtExpDay}{1556}         % us, median exposure of the day frames used
\newcommand{\sdRandBoxes}{338.2}           % us, random-position boxes, same shapes
\newcommand{\nullRandBoxes}{134.5}         % us, their analytic null
\newcommand{\excessRandBoxes}{306.3}       % us, their excess -- larger than the labels'
\newcommand{\framesRandBoxes}{1123}        % zurich_city_09_a.json, n_frames_with_rand
\newcommand{\minEvents}{200}               % events per box required
\newcommand{\minObjects}{3}                % qualifying objects per frame required

%% --- E09, the per-track structure of the frame-centered residual ------------
%% experiments/e09_per_object_evidence_time/per_track.py, re-run 2026-09-03 on
%% the corrected exposure window.  Computed after removing each frame's own
%% mean, so the frame-level component that a declared scalar timestamp would
%% cancel is removed by construction.
\newcommand{\evtTracks}{103}               % tracks with at least five observations
\newcommand{\betweenTrackSd}{149.2}        % us, between-track sd of the residual
\newcommand{\betweenTrackNull}{39.3}       % us, expected with no track effect
\newcommand{\trackVarRatio}{14.4}          % observed / no-effect variance ratio
\newcommand{\trackVarRatioCell}{8.6}       % same ratio after removing 32x32 cell means
\newcommand{\trackAcf}{0.721}              % lag-1 autocorrelation along a track
\newcommand{\acfLowDisp}{0.611}            % tracks below the displacement median
\newcommand{\acfHighDisp}{0.773}           % tracks above it
\newcommand{\dispSplit}{278}               % px, the median path length

%% --- E10, the mains line ---------------------------------------------------
\newcommand{\flickerPeriodUs}{10\,000}     % us, one 100 Hz period
\newcommand{\flickerPeakRatio}{10\,841}    % night peak over the 60-160 Hz continuum
\newcommand{\flickerLivePixNight}{35\,990} % live night pixels tested
\newcommand{\flickerLivePixDay}{17\,800}   % live day pixels tested
\newcommand{\flickerZpNinetyfive}{16.52}   % night 100 Hz p95 of Z
\newcommand{\flickerNullPninetyfive}{3.00} % Exp(1) p95
\newcommand{\flickerNullRate}{0.001}       % Exp(1) rate at p<1e-3
\newcommand{\flickerFracLockedSix}{0.0750} % fraction of live night pixels, p<1e-6
\newcommand{\flickerFracLockedDay}{0.0008} % day 100 Hz, same test

%% --- E11, the three pixel sets (result_exact.json) -------------------------
\newcommand{\excessMaskPix}{14\,334}       % phase-locked pixels excluded
\newcommand{\sdAll}{206.6}
\newcommand{\nullAll}{93.0}
\newcommand{\fracAll}{0.9012}
\newcommand{\sdClean}{205.1}
\newcommand{\nullClean}{97.1}
\newcommand{\fracClean}{0.8942}
\newcommand{\sdRandCtl}{206.3}
\newcommand{\nullRandCtl}{94.8}
\newcommand{\fracRandCtl}{0.8967}
\newcommand{\framesEleven}{1134}           % ALL and CLEAN arms
\newcommand{\framesRandCtl}{1133}          % matched random exclusion

%% --- E12, four analysis-window widths, RETRACTED with the block above -------
%% \newcommand{\wHalfPeriod}{5000}      %% \sdHalfPeriod 218.1  \nullHalfPeriod 46.5
%% \newcommand{\wOnePeriod}{10\,000}    %% \sdOnePeriod  211.5  \nullOnePeriod  74.4
%% \newcommand{\wFullExp}{14\,996}      %% \sdFullExp    206.6  \nullFullExp    93.0
%% \newcommand{\wTwoPeriod}{20\,000}    %% \sdTwoPeriod  271.2  \nullTwoPeriod 108.0
%% frames 1085 / 1118 / 1134 / 1137, frac over null 0.988 / 0.961 / 0.901 / 0.945

%% --- E08, object speed on DSEC-Det -----------------------------------------
\newcommand{\labelSpeedMedian}{40}         % px/s, median box-center speed
\newcommand{\labelSpeedPninety}{150}       % px/s, p90
\newcommand{\labelSpeedPninetynine}{360}   % px/s, p99
\newcommand{\offsetAtMedian}{1.0}          % px, 25 ms at the median speed
\newcommand{\offsetAtPninetynine}{9.0}     % px, 25 ms at p99

%% --- E13, the label center noise, solved rather than bounded ---------------
%% Difference-order ladder plus a two-rate elimination, the ladder is still
%% decreasing at k=7, so the total is a geometric extrapolation.
\newcommand{\sigmaCtotal}{0.4974}          % px, total center noise
\newcommand{\sigmaCquant}{0.1443}          % px, the integer-lattice component
\newcommand{\sigmaCannot}{0.476}           % px, the annotation component
\newcommand{\sigmaClattice}{0.5}           % px, the lattice step DSEC-Det stores on
\newcommand{\sigmaCmillis}{11.9}           % ms, sigmaCannot at the median speed
\newcommand{\ladderK}{7}                   % highest difference order taken

%% --- E14, the dispersion on the pixel scale --------------------------------
%% One multiplication: displacement = dispersion x object speed.
\newcommand{\pxAtMedian}{0.0073}           % px, at the median object speed
\newcommand{\pxAtPninety}{0.0275}          % px, at p90
\newcommand{\pxAtPninetynine}{0.0661}      % px, at p99
\newcommand{\pxFracMedian}{1.5}            % percent of sigmaCtotal
\newcommand{\pxFracPninetynine}{13.3}      % percent of sigmaCtotal
\newcommand{\speedToQuant}{786}            % px/s, speed at which it reaches sigmaCquant
\newcommand{\speedToNoise}{2709}           % px/s, speed at which it reaches sigmaCtotal

%% --- E16, the common offset from mid-exposure ------------------------------
\newcommand{\offsetAllEvents}{\ensuremath{-16.8}}    % us, night, all events
\newcommand{\offsetInBoxes}{\ensuremath{-102.3}}     % us, night, inside labeled boxes
\newcommand{\offsetDayBound}{2}            % us, both day figures fall below this
\newcommand{\offsetFramesSixteen}{1811}    % night frames surveyed
\newcommand{\pxOffsetPninetynine}{0.037}   % px, offsetInBoxes at p99 speed

%% --- macros for quantities that were still inline numerals ------------------
\newcommand{\dsecSeqCount}{eighteen}       % E00, sequences whose exposure file resolves
\newcommand{\dsecFrameRate}{20}            % Hz, label and frame rate
\newcommand{\ddArchiveBytes}{4\,655\,023}  % bytes, dsec-det_left_object_detections.zip
\newcommand{\ddArchiveDate}{2026-09-01}    % retrieval date
\newcommand{\sigmaEvtFramesLabelled}{1557} % E09, labelled frames in the night sequence
\newcommand{\mainsHz}{50}                  % Hz, Swiss mains
\newcommand{\flickerBandLo}{90}            % Hz, search band for the strongest line
\newcommand{\flickerBandHi}{110}
\newcommand{\flickerContLo}{60}            % Hz, continuum band
\newcommand{\flickerContHi}{160}
\newcommand{\flickerPeakRatioDay}{6.6}     % day peak over the same continuum
\newcommand{\flickerMinEvents}{100}        % events per pixel required
\newcommand{\flickerCtrlNightRate}{0.0004} % off-frequency arm, rate at p<1e-3
\newcommand{\periodsFullExp}{1.4996}   % E10, mains periods inside the ceiling exposure
\newcommand{\rvtCentroidMag}{25}           % ms, |rvtUniformCentroid|
\newcommand{\rvtTwoBranchScale}{7.2}       % ms, rvtCentroidMag/sqrt(12)

%% ===== E17 measured 2026-09-03: the network's temporal influence profile ====
%% 480 samples of real gen1 val, released rvt-t, recurrent state warmed eight
%% steps, each of the ten bins occluded in turn. This converts E02's
%% uniform-weight REFERENCE POINT into a MEASUREMENT.
\newcommand{\rvtCentroidMeas}{\ensuremath{-23.81}}   %% ms, influence-weighted centroid,
                                           %% E45. E17 and E34 reported -24.94 because the
                                           %% occluded pass received the recurrent state the
                                           %% reference pass RETURNED, not the one entering
                                           %% the step, E45 repeats them with the entering
                                           %% state and supersedes both.
\newcommand{\rvtCentroidMagMeas}{23.81}    %% ms, its magnitude, for "before"
\newcommand{\rvtCentroidUnif}{\ensuremath{-25.00}}   %% ms, uniform-weight reference
\newcommand{\rvtCentroidDiff}{\ensuremath{+1.19}}    %% ms, measured minus uniform
\newcommand{\rvtCentroidDiffPct}{23.8}     %% percent of one 5 ms bin: 1.190/5.0
\newcommand{\rvtInfluenceSamples}{480}
\newcommand{\rvtWarmSteps}{eight}          %% recurrent warm-up steps
\newcommand{\rvtInfluenceCV}{0.131}   %% coefficient of variation across the ten bins
\newcommand{\rvtInfluenceMaxMin}{1.52} %% heaviest bin over lightest
\newcommand{\rvtHalfOlder}{48.6}       %% percent of influence in the five older bins
\newcommand{\rvtHalfNewer}{51.4}       %% percent in the five newer bins
\newcommand{\tauHatRVT}{\rvtCentroidMeas}  %% E17, no longer a placeholder

%% --- proportions, written as percentages everywhere in the paper -----------
%% Same measurements as the fractions above, one format is used throughout so
%% that adjacent proportions are comparable by eye.
\newcommand{\sigmaEvtFracNightPct}{90.1}   % E11 exact, ALL arm
\newcommand{\sigmaEvtFracDayPct}{35}       % E09 day arm
\newcommand{\fracAllPct}{90.1}
\newcommand{\fracCleanPct}{89.4}
\newcommand{\fracRandCtlPct}{89.7}
\newcommand{\flickerFracLockedPct}{24.07}      % E10, of LIVE pixels
\newcommand{\flickerFracLockedDayPct}{0.08}    % E10, day arm
\newcommand{\flickerFracLockedSixPct}{7.50}    % E10, p<1e-6
\newcommand{\flickerCtrlNightRatePct}{0.04}    % E10, off-frequency arm
\newcommand{\flickerNullRatePct}{0.1}          % Exp(1) rate at p<1e-3
\newcommand{\cellGrid}{32\,$\times$\,32}       % E09, image-cell granularity

%% --- the ratio of the two temporal terms (E17 over E11) -------------------
\newcommand{\termRatio}{130}               % 23.81 ms / 183.6 us
\newcommand{\pxPredMedian}{1.0}            % px, tauHatRVT at the median speed
\newcommand{\pxPredPninetynine}{9.0}       % px, tauHatRVT at p99

%% ===== E19 measured 2026-09-03: the scalar dispersion against each box's own
%% mains modulation (experiments/e19_modulation_strata). For each box the
%% modulation depth at 100 Hz is m = 2|C| with C = (1/N) sum exp(-i w t_k) over
%% its own events, the within-frame dispersion is then recomputed inside
%% quartiles of m. This is what demotes the scalar dispersion to a bounded
%% statement: it rises monotonically with m and no unmodulated stratum exists.
\newcommand{\modFrames}{1134}              %% E19, frames entering the strata
\newcommand{\modBoxes}{6841}               %% E19, boxes entering the strata
\newcommand{\modPten}{0.355}               %% E19, p10 of m at 100 Hz
\newcommand{\modMedian}{0.467}             %% E19, median of m at 100 Hz
\newcommand{\modPninety}{0.609}            %% E19, p90 of m at 100 Hz
\newcommand{\modFloor}{0.125}              %% E19, m for 200 events at random phase
\newcommand{\excessModQone}{93.1}          %% us, E19, lowest 100 Hz quartile
\newcommand{\excessModQtwo}{128.1}         %% us, E19, second quartile
\newcommand{\excessModQthree}{131.5}       %% us, E19, third quartile
\newcommand{\excessModQfour}{208.0}        %% us, E19, highest quartile
\newcommand{\sdModQone}{119.3}             %% us, E19, sd of the lowest quartile
\newcommand{\sdModQtwo}{149.1}
\newcommand{\sdModQthree}{156.6}
\newcommand{\sdModQfour}{227.2}
\newcommand{\nullModQone}{74.6}            %% us, E19, its analytic null
\newcommand{\nullModQtwo}{76.2}
\newcommand{\nullModQthree}{85.1}
\newcommand{\nullModQfour}{91.5}
\newcommand{\framesModQone}{270}           %% E19, frames per stratum
\newcommand{\framesModQtwo}{243}
\newcommand{\framesModQthree}{207}
\newcommand{\framesModQfour}{265}
\newcommand{\modQoneLo}{0.122}             %% E19, the quartile edges of m
\newcommand{\modQoneHi}{0.410}
\newcommand{\modQtwoHi}{0.467}
\newcommand{\modQthreeHi}{0.534}
\newcommand{\modQfourHi}{1.050}
\newcommand{\modExcessRatio}{2.2}          %% E19, Q4 over Q1 at 100 Hz
\newcommand{\excessShamQone}{150.3}        %% us, E19, 137 Hz control, lowest quartile
\newcommand{\excessShamQfour}{187.6}       %% us, E19, 137 Hz control, highest quartile
\newcommand{\shamExcessRatio}{1.25}        %% E19, its Q4 over Q1
\newcommand{\flickerCtrlHz}{137}           %% Hz, the off-frequency of E10 and E19

%% ===== E20 RETRACTED 2026-09-05 ==============================================
%% The spatial gradient of evidence time. Five controlled passes, E25 generalized
%% the measurement to all six ceiling sequences and found the alignment STRONGER
%% where objects move slower, which a per-object motion gradient cannot produce,
%% and E26 compared each box against an equal-area annulus of its own surrounding
%% background: in four of six sequences the background aligns better than the
%% object, and eroding the box to its object-dominated core makes it worse in five
%% of six. The signal is ego-motion. Every macro of this block is left undefined so
%% that no sentence of the paper can use one. See experiments/e20_spatial_gradient
%% and docs/PROTOCOL_LEDGER.md entry 6.
%% \newcommand{\gradBoxes}{6627}              %% E20, qualifying boxes
%% \newcommand{\gradMinEvents}{400}           %% E20, events per box required
%% \newcommand{\gradSpeedFloor}{5}            %% px/s, E20, label-speed floor
%% \newcommand{\gradSpeedMedian}{58.3}        %% px/s, E20, median label speed
%% \newcommand{\gradCos}{\ensuremath{+0.152}} %% E20, mean cos(grad t, v)
%% \newcommand{\gradCosMedian}{\ensuremath{+0.304}}   %% E20, its median
%% \newcommand{\gradCosNullVel}{\ensuremath{+0.028}}  %% E20, shuffled-velocity null
%% \newcommand{\gradCosNullTime}{\ensuremath{+0.020}} %% E20, time-shuffled null
%% \newcommand{\gradSE}{0.012}                %% E20, 1/sqrt(n), a bound on the SE
%% \newcommand{\gradSigmas}{12.4}             %% E20, measured value in those units
%% \newcommand{\gradNullSigmas}{two}          %% E20, both nulls lie within this many
%% \newcommand{\gradMag}{11.23}               %% us/px, E20, median |grad t|
%% \newcommand{\gradMagShuffle}{5.47}         %% us/px, E20, its time-shuffled floor
%% \newcommand{\gradMagRatio}{2.05}           %% E20, measured over floor
%% \newcommand{\gradTranslatePx}{0.87}        %% px, gradSpeedMedian over dsecExpMax
%% \newcommand{\gradBoxSide}{42}              %% px, median sqrt(wh) of the released

%% --- the frame drawn in the qualitative figure -----------------------------
%% src/make_fig_qualitative.py, which ranks the sequence's candidate frames by
%% their within-frame dispersion and draws the MEDIAN one.
\newcommand{\figFiveCands}{602}            %% candidate frames ranked
\newcommand{\figFiveBoxes}{six}            %% qualifying boxes in the drawn frame
\newcommand{\figFiveEvents}{407\,942}      %% events inside its exposure window
\newcommand{\figFiveSd}{220.3}             %% us, its within-frame sd

%% ===== RETRACTED 2026-09-07: the E21, E23 and E24 blocks.
%% E24 fitted |e_par| and |e_perp|, E[|eps+tau v|] is not E[|eps|]+tau v, so its
%% coefficient difference does not identify a temporal offset (external review #3,
%% Critical 1, and the internal checklist audit of the same day). E23's acceleration
%% regressor was |dv|/dt, a magnitude entered where a signed a_par is required, and it
%% correlates with speed at +0.36. E21 and E24 both took the velocity from a FORWARD
%% difference, which shares the label-center noise with the residual and manufactures
%% a lag of tens of milliseconds when none is present (E29). All three are replaced by
%% E27/E30/E31/E32 below. The macros are removed rather than commented out so that no
%% text can reach a withdrawn number, see docs/REVIEW3_RESPONSE.md and
%% docs/PROTOCOL_LEDGER.md case 7.

%% ===== E27/E30/E31/E32 measured 2026-09-07: the effective OUTPUT time of the
%% released checkpoint, re-measured after three defects were found in E21 and E24.
%% Raw rows: experiments/e27_rows/rows.npz (35 523 matches, 406 sequences), one
%% GPU pass, every fit below is a CPU re-analysis of that one table.
%%   fit.log            the specification ladder and the sign-flip control
%%   placebo_diag.log   what the rotated-velocity coefficient is
%%   injection.log      recovery of an injected lag, and the hypothesis tests
%% Convention, stated once and used everywhere from here: p - g = -tau v, so a
%% POSITIVE tau means the output describes an EARLIER state than the label instant.
\newcommand{\outMatches}{22\,534}          %% E27, matched boxes with a centered velocity
\newcommand{\outRowsTotal}{35\,523}        %% E27, matches before requiring both neighbors
\newcommand{\outSeqs}{376}                 %% E27, sequences contributing to the fit
\newcommand{\outSeqsScanned}{406}          %% E27, sequences with at least one match
\newcommand{\outIoUfloor}{0.3}             %% E27, overlap required to match
\newcommand{\outIoUmedian}{0.873}          %% E27, median overlap of the matches
\newcommand{\outSpeedMedian}{5.4}          %% px/s, E27, median label speed, no floor
\newcommand{\outLinkPct}{79.0}             %% percent of Gen1 val boxes with a successor
                                           %% at IoU>=0.3, E36 (e27_rows/split_stats.log).
                                           %% The 81.5 here before was E21's floored subset.
\newcommand{\outLinkIoU}{0.3}              %% the overlap at which they link
\newcommand{\outSpeedSizeCorr}{\ensuremath{+0.139}} %% corr(|v|,sqrt(A)) on the E27 fit
                                           %% sample, E36. E21's +0.349 came from a
                                           %% different sample and does not apply here.

%% the specification ladder, centered velocity throughout (fit.log, placebo_diag.log)
\newcommand{\ladderMinTau}{\ensuremath{+14.62}}   \newcommand{\ladderMinTauSE}{5.47}
\newcommand{\ladderMinPl}{\ensuremath{-7.75}}     \newcommand{\ladderMinPlSE}{1.91}
\newcommand{\ladderMinPlSigmas}{4.0}
\newcommand{\ladderGeoTau}{\ensuremath{-7.62}}    \newcommand{\ladderGeoTauSE}{7.01}
\newcommand{\ladderGeoPl}{\ensuremath{-4.12}}     \newcommand{\ladderGeoPlSE}{1.90}
\newcommand{\ladderFullTau}{\ensuremath{-2.40}}   \newcommand{\ladderFullTauSE}{7.18}
\newcommand{\ladderFullPl}{\ensuremath{-2.66}}    \newcommand{\ladderFullPlSE}{1.71}
\newcommand{\ladderFullPlSigmas}{1.6}
\newcommand{\outTauSigmas}{0.33}           %% E31, tau in standard errors from zero
\newcommand{\outCentroidSigmas}{3.65}      %% E31, standard errors from tau = 23.81 ms
\newcommand{\outInterval}{26.2}            %% ms, evidence centroid minus output time
\newcommand{\outIntervalSE}{7.2}           %% ms, its standard error

%% the sign-flip control: the same fit under three velocity estimators (fit.log arm 1,
%% injection.log arm 3). A forward difference shares the label-center noise with the
%% residual, a backward one carries it with the opposite sign, a centered one omits it.
\newcommand{\sfMinFwd}{\ensuremath{-14.68}}   \newcommand{\sfMinFwdSE}{4.99}
\newcommand{\sfMinBack}{\ensuremath{+45.65}}  \newcommand{\sfMinBackSE}{5.59}
\newcommand{\sfMinCen}{\ensuremath{+17.51}}   \newcommand{\sfMinCenSE}{5.06}
\newcommand{\sfMinPlFwd}{\ensuremath{-6.53}} \newcommand{\sfMinPlFwdSE}{1.68}
\newcommand{\sfMinPlBack}{\ensuremath{-6.15}}\newcommand{\sfMinPlBackSE}{1.71}
\newcommand{\sfMinPlCen}{\ensuremath{-7.76}} \newcommand{\sfMinPlCenSE}{1.92}
\newcommand{\sfFullFwd}{\ensuremath{-20.60}}  \newcommand{\sfFullFwdSE}{6.36}
\newcommand{\sfFullBack}{\ensuremath{+14.38}} \newcommand{\sfFullBackSE}{5.97}
\newcommand{\sfFullPlFwd}{\ensuremath{-1.99}} \newcommand{\sfFullPlFwdSE}{1.39}
\newcommand{\sfFullPlBack}{\ensuremath{-2.10}}\newcommand{\sfFullPlBackSE}{1.50}
\newcommand{\sfSpread}{35.0}               %% ms, forward to backward under full controls
\newcommand{\sfMidpoint}{\ensuremath{-3.11}}  %% ms, their midpoint, against -2.40 centered

%% the placebo diagnosis (placebo_diag.log): splitting by the sign of the horizontal
%% velocity flips the rotated-velocity coefficient, which a rotational term cannot do.
\newcommand{\plRight}{\ensuremath{-27.99}} \newcommand{\plRightSE}{3.82}
\newcommand{\plRightSigmas}{7.3}
\newcommand{\plLeft}{\ensuremath{+21.84}}  \newcommand{\plLeftSE}{5.47}
\newcommand{\plLeftSigmas}{4.0}

%% injection (injection.log arms 1 and 2): the controls do not absorb a lag.
\newcommand{\injRecovery}{1.000}           %% recovered fraction at 10, 25, 50, 100 ms
\newcommand{\injGeoGamma}{0.05}            %% px per px, the geometric displacement injected
\newcommand{\injGeoMin}{\ensuremath{+0.63}}%% ms, how far it moves the minimal specification
\newcommand{\injGeoFull}{0.01}             %% ms, how far it moves the full one, below this

%% the speed-floor arm (placebo_diag.log section 6), full controls throughout
\newcommand{\floorFiveTau}{\ensuremath{-2.92}}    \newcommand{\floorFiveTauSE}{9.05}
\newcommand{\floorTwentyTau}{\ensuremath{-20.24}} \newcommand{\floorTwentyTauSE}{14.78}

%% ===== E28 measured 2026-09-07 (experiments/e28_gen1_label_noise): the label-center
%% noise of Gen1, measured on Gen1 by E13's difference ladder. E13's 0.4974 px is
%% DSEC-Det's and does not transfer.
\newcommand{\labelSigmaGen}{0.6515}        %% px per axis, E28, sigma_total
\newcommand{\labelSigmaGenAnn}{0.6353}     %% px, its annotation part
\newcommand{\labelSigmaGenQuant}{0.1443}   %% px, its quantization part, 0.5 px lattice
\newcommand{\labelTracksGen}{691}          %% E28, tracks of >= 12 labels on uniform 250/500/1000-ms grids

%% ===== E29 measured 2026-09-07 (experiments/e29_synthetic_signflip): the
%% forward-difference artefact, calibrated on synthetic tracks with tau_true = 0.
\newcommand{\synFwd}{\ensuremath{-5.64}}   %% ms, forward estimator, no lag present
\newcommand{\synBack}{\ensuremath{+5.84}}  %% ms, backward estimator, no lag present
\newcommand{\synCen}{\ensuremath{+0.00}}   %% ms, centered estimator, no lag present
\newcommand{\synSum}{\ensuremath{+0.20}}   %% ms, forward plus backward
\newcommand{\synPred}{\ensuremath{-5.63}}  %% ms, -Var(delta)/dt / Var(v), closed form
\newcommand{\synZero}{\ensuremath{+0.02}}  %% ms, all three with the label noise removed
\newcommand{\synRecover}{0.92}             %% fraction of an injected lag recovered
\newcommand{\synLambda}{0.933}             %% the errors-in-variables prediction
\newcommand{\synTracks}{2000}              %% synthetic tracks

%% displacement the interval amounts to, at speeds that occur in the split (E32)
\newcommand{\outSideMedian}{41.7}          %% px, median box side of the matches
\newcommand{\outSpeedNinety}{25.7}         %% px/s, 90th percentile label speed
\newcommand{\outSpeedNinetynine}{54.2}     %% px/s, 99th percentile label speed
\newcommand{\outPxNinety}{0.61}            %% px, displacement over 23.81 ms at p90
\newcommand{\outPxNinetynine}{1.29}        %% px, the same at p99
\newcommand{\outStaticPct}{22}             %% percent of matches below 2 px/s

%% ===== E25 measured 2026-09-04 (experiments/e25_six_sequences): the mains
%% signature on ALL SIX DSEC train sequences whose exposure is pinned at 14996 us.
%% Per labelled box, the modulation depth m = 2|C| and the Rayleigh statistic
%% Z = N|C|^2 at 100 Hz, with an off-frequency arm at 137 Hz. For these finite
%% exposure windows the per-box analytic phase null is not used; E38 supplies
%% an empirical off-frequency reference.
\newcommand{\sixSeqs}{six}                 %% E25, ceiling sequences measured
\newcommand{\sixBoxes}{16\,516}            %% E25, labelled boxes over the six
\newcommand{\sixZmedian}{114.58}           %% E25, median over sequences of the median box Z at 100 Hz
\newcommand{\sixZctrl}{4.01}               %% E25, the same at 137 Hz
\newcommand{\sixLockedLoPct}{95.4}         %% percent, E38b: the LOWEST per-sequence fraction of boxes whose local excess
                                           %% at 100 Hz exceeds the 99th percentile of the
                                           %% same quantity at 137 Hz. This replaces E25's
                                           %% p<1e-3 fraction, withdrawn with the analytic
                                           %% per-box null (experiments/e38_empirical_null).
\newcommand{\sixLockedHiPct}{100.0}        %% percent, E38b, highest (same construction)
\newcommand{\sixModLo}{0.428}              %% E25, lowest per-sequence median modulation depth
\newcommand{\sixModHi}{0.489}              %% E25, highest
%% per-sequence rows of Table 2, in the order printed
\newcommand{\sixBoxesA}{2914}   \newcommand{\sixModA}{0.428}
\newcommand{\sixZA}{117.6}      \newcommand{\sixZcA}{3.15}  \newcommand{\sixLockA}{100.00}
\newcommand{\sixBoxesB}{2134}   \newcommand{\sixModB}{0.440}
\newcommand{\sixZB}{99.1}       \newcommand{\sixZcB}{4.69}  \newcommand{\sixLockB}{99.63}
\newcommand{\sixBoxesC}{619}    \newcommand{\sixModC}{0.445}
\newcommand{\sixZC}{68.4}       \newcommand{\sixZcC}{2.75}  \newcommand{\sixLockC}{99.68}
\newcommand{\sixBoxesD}{241}    \newcommand{\sixModD}{0.489}
\newcommand{\sixZD}{310.9}      \newcommand{\sixZcD}{25.85} \newcommand{\sixLockD}{95.44}
\newcommand{\sixBoxesE}{6997}   \newcommand{\sixModE}{0.466}
\newcommand{\sixZE}{251.5}      \newcommand{\sixZcE}{7.82}  \newcommand{\sixLockE}{99.23}
\newcommand{\sixBoxesF}{3611}   \newcommand{\sixModF}{0.438}
\newcommand{\sixZF}{111.6}      \newcommand{\sixZcF}{3.33}  \newcommand{\sixLockF}{100.00}
\newcommand{\sixMinEvents}{400}            %% E25, events per box required

%% ===== E26 measured 2026-09-04 (experiments/e26_egomotion, paired.json): the
%% control that retracted E20. Each box's fitted gradient is compared against the
%% gradient of an equal-area annulus of its own surrounding background, aligned
%% with the same label velocity, and the difference is taken per box so that a
%% cause common to the box and its surroundings cancels.
\newcommand{\egoRingBetter}{four}          %% E26, sequences where the ring aligns better
\newcommand{\egoCoreNeg}{five}             %% E26, sequences where the eroded core goes negative
\newcommand{\egoCoreWorst}{\ensuremath{-0.178}}  %% E26, paired core minus ring, zurich_city_01_a
\newcommand{\egoCoreWorstSigmas}{6.2}      %% E26, that difference in standard errors

%% --- E02, provenance of the window facts (src/e02_verify_window.py) ---------
\newcommand{\rvtSourceLine}{405}           %% line of preprocess_dataset.py that builds
                                           %% ev_repr_timestamps_us_end backwards from labels

%% --- the two temporal terms on the pixel scale of the split they were measured on

%% ===== E32b measured 2026-09-07 (experiments/e27_rows/bootstrap.log): a cluster
%% bootstrap over sequences for the reported specification, fitted by
%% Frisch-Waugh-Lovell (the group dummies are absorbed by within-cell demeaning,
%% which reproduces the full design's tau to 0.007 ms).
\newcommand{\outBootLo}{\ensuremath{-17.0}}   %% ms, 2.5th percentile
\newcommand{\outBootHi}{\ensuremath{+11.5}}   %% ms, 97.5th percentile
\newcommand{\outBootDraws}{4000}              %% resamples
\newcommand{\outBootBeyond}{3}                %% draws at or beyond 23.81 ms, of 4000

%% ===== E33 measured 2026-09-07 (experiments/e27_rows/saturate.log): the reported
%% specification pairs each geometric control with one image axis. This enters all
%% four on BOTH axes, drops each in turn, and restricts the sample, so that the
%% spread of tau across specifications is reported rather than one of them chosen.
\newcommand{\satTau}{\ensuremath{-3.62}}   \newcommand{\satTauSE}{7.41}
\newcommand{\satPl}{\ensuremath{-1.59}}    \newcommand{\satPlSE}{1.62}
\newcommand{\satLo}{\ensuremath{-20.2}}    %% ms, lowest tau over EVERY specification
                                           %% carrying geometric controls, including the
                                           %% 20 px/s floor of placebo_diag.log, the
                                           %% saturate.log sweep alone gives -9.8.
\newcommand{\satHi}{\ensuremath{+11.1}}    %% ms, highest
\newcommand{\satTight}{\ensuremath{+0.49}} \newcommand{\satTightSE}{4.67}
\newcommand{\satTightIoU}{0.7}             %% the overlap floor giving the tightest fit

%% ===== E34 measured 2026-09-07 (experiments/e34_influence_robust/result.json): is the
%% influence centroid an artefact of occluding a bin with zeros? Three fills and one
%% instrument that does not occlude at all. The recurrent state is held fixed within
%% every sample, so none of these measures state perturbation.
\newcommand{\occZero}{\ensuremath{-23.81}}   %% ms, zeros, what E17 did
\newcommand{\occMean}{\ensuremath{-26.35}}   %% ms, the bin set to its dataset mean
\newcommand{\occSwap}{\ensuremath{-25.01}}   %% ms, the bin taken from an unrelated sample
\newcommand{\occGrad}{\ensuremath{-22.35}}   %% ms, d||out||/d(bin), no occlusion at all
\newcommand{\occSpread}{4.01}                 %% ms, the range across all four instruments
\newcommand{\occBootSE}{0.054}                %% ms, bootstrap SE of the zero-fill centroid
\newcommand{\occGradCV}{0.191}               %% CV over bins of the gradient instrument

%% ===== E35 measured 2026-09-07 (experiments/e27_rows/ratio.log): the nearest prior
%% estimator, the per-object ratio, fitted on exactly the rows the regression uses.
\newcommand{\ratioMean}{\ensuremath{+24.22}} \newcommand{\ratioMeanSE}{27.73}
\newcommand{\ratioIQR}{430}                  %% ms, its interquartile range over all rows
\newcommand{\ratioSD}{4170}                  %% ms, its sample standard deviation
\newcommand{\ratioFloor}{\ensuremath{+13.34}}\newcommand{\ratioFloorSE}{6.17}
\newcommand{\ratioFloorSpeed}{15}            %% px/s, the floor that arm imposes

%% plot and procedure constants, so that no numeral appears in prose
\newcommand{\bootCI}{95}                     %% percent, the bootstrap interval reported
\newcommand{\injLevels}{10, 25, 50 and 100}  %% ms, the injected lags of injection.log
\newcommand{\flickerDayPeakHz}{103.0}        %% Hz, where the day sequence's in-band peak sits
                                             %% (e10_flicker/verify.json, DAY top_line_Hz)

%% errors-in-variables on the reported rows (experiments/e27_rows/fit.log, section 5,
%% re-run with SIGMA_C from E28 rather than the E13 default)
\newcommand{\eivVarObs}{110.6}             %% px^2/s^2, Var of the centered velocity
\newcommand{\eivVarNoise}{1.0}             %% px^2/s^2, row-wise mean 2 sigma^2 / (dt_f+dt_b)^2
\newcommand{\eivLambda}{0.991}             %% the attenuation factor on these rows
\newcommand{\eivShift}{0.2}                %% ms, conservative rounded upper bound on correction shift

%% ===== E38 measured 2026-09-07 (experiments/e38_empirical_null): the per-box Rayleigh
%% null, measured rather than assumed. RETRACTION: the per-box p-values and the
%% "99.0-100.0 % of boxes locked at p<1e-3" figure are withdrawn. Pooled over the six
%% ceiling sequences the OFF-FREQUENCY per-box Z has a median of 38.39 against the
%% analytic 0.693, an inflation of 55, because event times inside one box during a 15 ms
%% exposure arrive in bursts. The per-PIXEL arm is unaffected: there the 137 Hz control
%% returns 0.581 against 0.693, so the analytic null holds where it was verified.
\newcommand{\nullOffMed}{38.4}             %% pooled per-box Z at off frequencies
\newcommand{\nullInflation}{55}            %% times the analytic median
\newcommand{\nullOffFreqs}{eight}             %% off-frequencies swept

%% arm 2 (local_excess.log): R = Z(100 Hz) / the same box's own local background, which
%% needs no null at all. Grid 60-300 Hz at 1 Hz, the background band is |f-100| in [6,30]
%% excluding 20 Hz and 100 Hz harmonics.
\newcommand{\localBoxes}{16\,516}          %% boxes pooled over the six sequences
\newcommand{\localLine}{1.87}              %% median R at the line
\newcommand{\localLinePten}{1.62}          %% its tenth percentile
\newcommand{\localLinePninety}{2.18}       %% its ninetieth
\newcommand{\localSham}{0.11}              %% median R at the off-frequency arm
\newcommand{\localShamPninetynine}{1.00}   %% its 99th percentile
\newcommand{\localRatio}{17.5}             %% ratio of the two medians
\newcommand{\localAbovePct}{99.63}         %% percent of boxes above the sham's 99th pct
\newcommand{\localPeakHz}{96}              %% Hz, where the pooled median spectrum peaks
\newcommand{\localPeakOverHigh}{8.5}       %% peak over the 200-300 Hz median
\newcommand{\localPeakWidth}{35}           %% Hz, its half-maximum width
\newcommand{\localTransformLimit}{67}      %% Hz, 1/14996 us, what one window can resolve
\newcommand{\localGridLo}{60}\newcommand{\localGridHi}{300}

%% ===== E37 measured 2026-09-07 (experiments/e37_map): average precision as a function
%% of the instant the ground truth describes. One GPU pass dumps every detection and every
%% ground-truth box of the Gen1 validation split, the sweep displaces the ground truth by
%% delta*v and recomputes mAP. Ground truth without a centered velocity is marked ignore.
\newcommand{\mapFrames}{20\,296}           %% frames evaluated
\newcommand{\mapDets}{86\,788}             %% detections dumped
\newcommand{\mapGT}{40\,698}               %% ground-truth boxes
\newcommand{\mapGTvelPct}{68.7}            %% percent carrying a centered velocity
\newcommand{\mapArgmax}{\ensuremath{-10}}  %% ms, where mAP is maximised
\newcommand{\mapZero}{0.3346}              %% mAP at the label instant
\newcommand{\mapCentroid}{0.3342}          %% mAP with the gt displaced to the centroid
\newcommand{\mapCost}{0.04}                %% mAP points between those two
\newcommand{\mapCostAbs}{0.0004}           %% absolute mAP change
\newcommand{\mapCostPct}{0.12}             %% and as a percentage
\newcommand{\mapSweepSpan}{0.9}            %% mAP points over the whole +-60 ms sweep
\newcommand{\mapSpeedMedian}{4.0}          %% px/s, median label speed of this set
\newcommand{\mapSideMedian}{41.5}          %% px, median box side
\newcommand{\mapPxMedian}{0.10}            %% px, 23.81 ms at that median speed
%% the speed strata (strata.log): span over +-50 ms, and n
\newcommand{\mapSpanSlow}{0.11}   \newcommand{\mapNslow}{20\,154}
\newcommand{\mapSpanMid}{0.33}    \newcommand{\mapNmid}{5\,356}
\newcommand{\mapSpanHigh}{0.78}   \newcommand{\mapNhigh}{2\,096}
\newcommand{\mapSpanTop}{0.37}    \newcommand{\mapNtop}{337}
\newcommand{\mapTopMap}{0.139}    %% the fastest stratum's mAP, against 0.286 below it
\newcommand{\mapHighMap}{0.286}

%% ===== E55 measured 2026-09-15 (experiments/e37_map/stratum_table.json): the speed
%% strata evaluated EXACTLY at the corrected centroid of -23.81 ms (E37c's per-stratum
%% costs sat at the superseded -24.94 ms and on the 5 ms grid), with each stratum's
%% share of the moving boxes, its median speed, and what 23.81 ms of that speed is in
%% pixels. Spans are E37c's, whose grid did not change. The all-moving row reproduces
%% at_centroid.json to inside 5e-5, asserted in the script.
\newcommand{\stShareSlow}{72.1}\newcommand{\stShareMid}{19.2}
\newcommand{\stShareHigh}{7.5}\newcommand{\stShareTop}{1.2}
\newcommand{\stSpdSlow}{2.1}\newcommand{\stSpdMid}{15.0}
\newcommand{\stSpdHigh}{32.3}\newcommand{\stSpdTop}{57.4}
\newcommand{\stPxSlow}{0.05}\newcommand{\stPxMid}{0.36}
\newcommand{\stPxHigh}{0.77}\newcommand{\stPxTop}{1.37}
\newcommand{\stCostSlow}{0.02}\newcommand{\stCostMid}{0.01}
\newcommand{\stCostHigh}{0.08}\newcommand{\stCostTop}{0.00}
\newcommand{\stSpanAll}{0.37}     %% the all-moving span on the -50..+30 ms grid, the
                                  %% 0.9 of \mapSweepSpan is the +-60 ms sweep of E37b
\newcommand{\stNmoving}{27\,943}  %% ground-truth boxes carrying a centered velocity

%% ===== Real-image figures rendered 2026-09-15 (src/make_fig_real.py). fig6: the frame is
%% chosen by rule (>=3 velocity boxes, >=2 at >=20 px/s, largest median speed) from E37's
%% own iteration order, so the drawn boxes are the ones the sweep scored, stats in
%% experiments/e37_map/fig6_stats.json. fig7: densest published ceiling exposure window vs
%% the same instants of the daytime control, stats in
%% experiments/e00_exposure_survey/fig7_stats.json.
\newcommand{\sweepRealFast}{45.2}   %% px/s, the fastest label on the drawn frame
\newcommand{\sweepRealPx}{1.08}     %% px, that speed over the 23.81 ms interval
\newcommand{\ceilVarNight}{3.0}     %% max/min of the 350 us rate, ceiling sequence
\newcommand{\ceilVarDay}{1.1}       %% the same on the daytime control

%% ===== E40 measured 2026-09-07 (experiments/e40_harmonics/result.json): which frequency
%% carries the signature, with no band chosen in advance. E10's "strongest line between 90
%% and 110 Hz" could only ever have found a line in that band. A 20-450 Hz sweep on each
%% whole recording shows the HARMONIC SERIES of a 100 Hz fundamental, strongest at the
%% second harmonic, with nothing at the odd multiples of 50 Hz and nothing in daylight.
\newcommand{\harmOne}{30}                  %% ceiling median, 100 Hz over its own continuum
\newcommand{\harmTwo}{125}                 %% the same at 200 Hz
\newcommand{\harmThree}{15}                %% at 300 Hz
\newcommand{\harmFifty}{1.8}               %% at 50 Hz: no line
\newcommand{\harmOneFifty}{4.5}            %% at 150 Hz: no line
\newcommand{\harmTwoFifty}{1.6}            %% at 250 Hz: no line
\newcommand{\harmDayOne}{0.5}              %% daytime sequence at 100 Hz
\newcommand{\harmDayTwo}{6.2}              %% daytime sequence at 200 Hz
\newcommand{\harmTwoMax}{2518}             %% the strongest, zurich_city_09_a at 200 Hz
\newcommand{\harmTwoMin}{14.7}             %% the weakest ceiling sequence at 200 Hz
\newcommand{\harmSecondStrongest}{four}    %% sequences where 200 Hz is the strongest line
\newcommand{\harmGridLo}{20}\newcommand{\harmGridHi}{450}
\newcommand{\harmSecondHz}{200}\newcommand{\harmThirdHz}{300}
\newcommand{\harmBoxTwo}{0.76}             %% E38b at 200 Hz: the per-box local excess there
                                           %% is WEAKER than at 100 Hz (1.87), because one
                                           %% 14996 us window resolves only 67 Hz

%% ===== E41 measured 2026-09-07 (experiments/e41_permutation_null): the within-frame
%% dispersion against a null that assumes nothing about arrival statistics. Events of all
%% qualifying boxes of a frame are pooled and reassigned at random to boxes, preserving
%% counts, so every temporal structure common to the frame survives and only box identity
%% is destroyed. The analytic null's independence assumption is thereby removed.
\newcommand{\permNull}{78.6}               %% us, the permutation null, median over frames
\newcommand{\permExcess}{191.1}            %% us, the excess over it
\newcommand{\permDraws}{200}               %% draws per frame
\newcommand{\permFrames}{1134}             %% qualifying frames
\newcommand{\permAbovePct}{94.1}           %% percent of frames exceeding their own null
\newcommand{\permRatio}{0.85}              %% permutation over analytic: it is SMALLER
\newcommand{\permRandObs}{339.4}           %% us, random-position boxes, observed
\newcommand{\permRandNull}{113.3}          %% us, their permutation null
\newcommand{\permRandExcess}{320.0}        %% us, their excess

%% retained for provenance only; the finite-window per-box table reports depth
%% descriptively and does not use a random-phase floor.
\newcommand{\modFloorSix}{0.089}           %% sqrt(pi/400)
\newcommand{\fwlAgree}{0.007}              %% ms, FWL against the full design (bootstrap.log)
\newcommand{\mapArgmaxTau}{\ensuremath{+10}}  %% ms, \mapArgmax read on the p-g=-tau v axis
\newcommand{\localRatioSweep}{4.0}         %% the pooled 100 Hz to off-frequency ratio of
                                           %% medians over all eight off frequencies (E38
                                           %% run.log), against 28.6 for the 137 Hz arm alone
\newcommand{\mapVsFit}{1.7}                %% the mAP argmax against the regression, in SEs
                                           %% of the regression: (10-(-2.40))/7.18 = 1.73

%% ===== E42 measured 2026-09-07 (experiments/e42_recurrent_support): the temporal support
%% of a RECURRENT detector, including its history. E17 and E34 hold the recurrent state
%% fixed and perturb the current window only, so their centroid is the CURRENT WINDOW's
%% influence centroid conditional on that state, not the support of the emitted state as
%% Sec. 3.1 defines it. This occludes whole past windows and lets the state propagate.
\newcommand{\histShareTen}{30.0}           %% percent of influence in the current window,
                                           %% 500 ms horizon
\newcommand{\histShareTwenty}{22.7}        %% the same over a 1000 ms horizon
\newcommand{\histCentroidTen}{\ensuremath{-161}}    %% ms, over the FIRST TEN windows of the k=19
                                           %% run, so both horizons come from one sample,
                                           %% the k=9 run began at a different frame and its
                                           %% -166 is not mixed in
\newcommand{\histCentroidTwenty}{\ensuremath{-299}} %% ms, the same over 1000 ms
\newcommand{\histHalfLag}{\ensuremath{-175}}        %% ms, where half the influence is reached
\newcommand{\histResetChange}{0.0867}      %% relative change when the state is zeroed
\newcommand{\histCurrentChange}{0.0231}    %% the same for occluding the whole current window
\newcommand{\histResetRatio}{3.8}          %% the ratio of those two
\newcommand{\histSamples}{576}             %% samples
\newcommand{\histHorizonA}{500}\newcommand{\histHorizonB}{1000}  %% ms
%% the fill control (fill_control.log): the plateau is not an artefact of zero filling
\newcommand{\histDecayZero}{10.3}          %% decay from lag 0 to lag 19, zero fill
\newcommand{\histDecaySwap}{2.2}           %% the same under a swap fill
\newcommand{\histSwapRatioLo}{2.43}       %% swap/zero sensitivity ratio at lag 0
\newcommand{\histSwapRatioHi}{11.47}      %% the same at lag 19
\newcommand{\histTailSwapPct}{46}          %% percent of lag-0 influence still at lag 19

%% stratum edges and band frequencies, so that no numeral enters prose (E37 strata.log,
%% E40 result.json). The p<1e-3 level is a chosen significance and is written as a power.
\newcommand{\mapCutLo}{10}\newcommand{\mapCutMid}{25}\newcommand{\mapCutHi}{50}
\newcommand{\harmOddOneHz}{150}\newcommand{\harmOddTwoHz}{250}
\newcommand{\ladderTrimPct}{5}             %% percent trimmed per tail in E28's ladder

%% ===== E37d measured 2026-09-07 (experiments/e37_map/argmax_bootstrap.log): a frame
%% bootstrap of the mAP sweep's argmax. The argmax does not identify a value, and in no
%% resample does it reach the newest window's centroid.
\newcommand{\mapArgLo}{\ensuremath{-20}}   %% ms, 2.5th percentile of the argmax
\newcommand{\mapArgHi}{\ensuremath{0}}     %% ms, 97.5th percentile
\newcommand{\mapArgBoot}{1000}             %% frame resamples
\newcommand{\mapArgNever}{none}            %% resamples whose argmax reaches -23.81 ms

%% ===== E44 measured 2026-09-07 (experiments/e44_capacity/rvt-{t,s,b}.json): the same
%% predictor-side measurement on the three released Gen1 capacities, which differ only in
%% embed_dim, attention dim_head and FPN depth. Same sequences, representation, warm-up and
%% code for all three, all three load with 0 missing and 0 unexpected keys (E43). The
%% newest-window centroids here use the state entering the step, as E45 does, which is why
%% rvt-t agrees with E45 and not with the superseded E17.
\newcommand{\capSamples}{576}            %% samples per capacity
\newcommand{\capParamsT}{4.41}        \newcommand{\capParamsS}{9.87}
\newcommand{\capParamsB}{18.54}
\newcommand{\capBinT}{\ensuremath{-23.84}}   \newcommand{\capBinS}{\ensuremath{-23.89}}
\newcommand{\capBinB}{\ensuremath{-24.11}}
\newcommand{\capBinSpread}{0.27}   %% ms, range of the three
\newcommand{\capShareT}{22.7}   \newcommand{\capShareS}{20.6}
\newcommand{\capShareB}{19.0}
\newcommand{\capSuppT}{\ensuremath{-299}}  \newcommand{\capSuppS}{\ensuremath{-337}}
\newcommand{\capSuppB}{\ensuremath{-339}}
\newcommand{\capResetT}{3.8}   \newcommand{\capResetB}{3.0}
\newcommand{\capHorizon}{1000}          %% ms, the horizon these are quoted over

%% ===== E47/E48 measured 2026-09-08 (experiments/e47_ssm, experiments/e48_matched_frames):
%% the newest-window bin arm on TWO architectures over the same 12 Gen1 validation
%% sequences and the same frames. SSM-ViT (Zubic et al., CVPR 2024) is a fork of the RVT
%% repository in which each stage's ConvLSTM is replaced by an S5 state space layer, it
%% distributes RVT's own preprocessed Gen1 and builds its window boundaries with the same
%% line of the same script. Both released checkpoints load with 0 missing and 0 unexpected
%% keys. E48 re-runs RVT's three capacities on E47's frames so the two families are
%% measured on identical samples, its rvt-t agrees with E45. SSM-ViT is driven the way its
%% released evaluation drives it, a chunk of windows at a time with the detection head
%% applied per position, and measured at positions carrying half a chunk of history or more
%% (E47b), the one-window-at-a-time convention is withdrawn, see withdrawn_L1/WHY.md.
\newcommand{\archCkpts}{5}            %% released checkpoints measured
\newcommand{\archOther}{4}            %% the ones Fig. 2a draws besides rvt-t
\newcommand{\archSpan}{0.96}       %% ms, range of all their centroids
\newcommand{\archLo}{\ensuremath{-24.72}}   \newcommand{\archHi}{\ensuremath{-23.76}}
\newcommand{\ssmParamsS}{9.68}  \newcommand{\ssmParamsB}{18.19}
\newcommand{\ssmBinS}{\ensuremath{-23.76}}
\newcommand{\ssmBinB}{\ensuremath{-24.72}}
\newcommand{\mfBinT}{\ensuremath{-23.81}}
\newcommand{\mfBinS}{\ensuremath{-23.98}}
\newcommand{\mfBinB}{\ensuremath{-24.13}}
\newcommand{\archSamples}{480}          %% samples per RVT checkpoint
\newcommand{\archSamplesSsm}{288}       %% samples per SSM-ViT checkpoint:
                                           %% fewer, because only chunk positions carrying
                                           %% half a chunk of history or more are used
\newcommand{\archUnifMax}{1.24}       %% ms, the largest distance of any
                                           %% of them from the uniform-weight -25.00
\newcommand{\archUnifPct}{25}        %% that, as a percent of one bin
%% ===== E47c measured 2026-09-08 (experiments/e47_ssm/statepos-*.json): what the
%% recurrent state carried between evaluation chunks does to the output, by position
%% in the chunk, as a relative L2 change in the emitted detection tensor. The control
%% arm occludes the window one position earlier inside the same chunk, which a model
%% that ignored its own history entirely would also return zero for.
\newcommand{\ssmStateEffect}{0.00000}   %% largest over all positions, s5vit-small
\newcommand{\ssmInchunkOne}{0.0276}      %% previous window, position 1
\newcommand{\ssmInchunkLast}{0.0106}     %% the same at the last position
\newcommand{\ssmChunkMeas}{8}          %% windows per chunk in this measurement
\newcommand{\rvtStateEffect}{0.0304}     %% the same arm on rvt-t, last position
\newcommand{\rvtStateEffectOne}{0.0616}  %% and at position 1
\newcommand{\ssmEvalChunk}{21}          %% windows per chunk in the released streaming
                                           %% evaluation: config/dataset/gen1.yaml
                                           %% sequence_length, and sequence_for_streaming.py
                                           %% start_indices step by exactly that

%% ===== E51-E53 measured 2026-09-08 (experiments/e51_ranking): does the interval change a
%% benchmark COMPARISON? Detections for all five released checkpoints were dumped under one
%% protocol (same frames, same confidence, same NMS, ground truth identical across models and
%% identical to E37's), then each model was scored at delta = 0 and again at its own best
%% displacement, in four speed strata. The ordering never changes, and it cannot: in every
%% stratum the largest gain any model draws from alignment is a fraction of the smallest gap
%% between adjacent models, so no refinement of the grid can reorder them.
\newcommand{\rankCkpts}{5}            %% released checkpoints scored
\newcommand{\rankStrata}{four}          %% speed strata examined
\newcommand{\rankFlips}{0}              %% pairwise sign flips, of 10 per stratum
\newcommand{\rankMaxGain}{0.14}      %% mAP points, the largest alignment gain seen
\newcommand{\rankMinGap}{0.11}       %% mAP points, the smallest adjacent gap seen
\newcommand{\rankRatio}{52}          %% percent: worst-case gain over gap within a stratum
\newcommand{\rankArgLo}{\ensuremath{-15}}   %% ms, earliest per-model argmax
\newcommand{\rankArgHi}{\ensuremath{0}}    %% ms, latest per-model argmax
\newcommand{\rankMover}{\texttt{s5vit-base}}  %% the checkpoint whose rank moves most
\newcommand{\rankMoverBest}{1}          %% its best rank across strata
\newcommand{\rankMoverWorst}{5}         %% and its worst

%% ---------------------------------------------------------------------------
%% Resolving power of the metric (E56). What size of timing difference the mAP
%% can separate from its own cluster-resampling noise, over 406 sequences.
%% Emitted by src/e61_macros.py from experiments/e56_resolving/resolving.json.
\newcommand{\seDmapLo}{0.035}
\newcommand{\seDmapHi}{0.058}
\newcommand{\seMapLo}{2.19}
\newcommand{\seMapHi}{2.82}
\newcommand{\resPairedLo}{13.7}
\newcommand{\resPairedHi}{19.8}
\newcommand{\effTauLo}{0.041}
\newcommand{\effTauHi}{0.081}
\newcommand{\twoSeDmapLo}{0.070}
\newcommand{\twoSeDmapHi}{0.116}
\newcommand{\effTauResolved}{1}
\newcommand{\resAbsLo}{153}
\newcommand{\resAbsHi}{174}
\newcommand{\dstarAll}{43.7}
\newcommand{\dstarAllTau}{1.83}
\newcommand{\dstarMid}{11.2}
\newcommand{\dstarMidTau}{0.47}

%% Stability of the published ordering under the benchmark's own noise (E56d),
%% from experiments/e56_resolving/rankse.json.
\newcommand{\rankHoldBox}{42.5}
\newcommand{\rankHoldMov}{80.5}
\newcommand{\rankFlipMax}{0.48}
\newcommand{\rankGapMin}{0.050}
\newcommand{\rankGapMinSE}{0.53}

%% Chunk position in the released SSM-ViT streaming evaluation (E58), from
%% experiments/e58_chunkpos/{allbox,invariant}.json.
\newcommand{\chunkDidBase}{5.32}
\newcommand{\chunkDidBaseSE}{0.53}
\newcommand{\chunkDidBaseZ}{10.0}
\newcommand{\chunkDidSmall}{4.37}
\newcommand{\chunkDidSmallSE}{0.61}
\newcommand{\chunkDidSmallZ}{7.2}
%% Raw full-minus-short mAP difference per checkpoint (E58, all labeled boxes), the
%% quantity main Fig.~2 plots, kept separate from the control-adjusted difference below.
\newcommand{\chunkRawT}{0.01}
\newcommand{\chunkRawS}{0.80}
\newcommand{\chunkRawB}{0.56}
%% The same raw difference net of the placebo trend: for an RVT row the mean of the other
%% two RVT rows, for an SSM row the mean of all three. Positive means history helped.
\newcommand{\chunkNetT}{-0.67}
\newcommand{\chunkNetS}{0.51}
\newcommand{\chunkNetB}{0.16}
\newcommand{\chunkPlaceboMax}{0.67}
\newcommand{\chunkPlaceboZ}{1.56}
\newcommand{\chunkGapBase}{1.75}
\newcommand{\chunkGapSmall}{1.33}
\newcommand{\chunkShortBase}{34.445}
\newcommand{\chunkFullBase}{40.226}
\newcommand{\chunkRelBase}{38.9}
\newcommand{\chunkConfOne}{0.463}
\newcommand{\chunkConfFull}{0.574}
\newcommand{\chunkConfPct}{23.8}

%% Which summary of the label composition sets the resolution, and what buying it
%% costs (E56e), from experiments/e56_resolving/collapse.json and composition.json.
\newcommand{\collapseCVraw}{0.20}
\newcommand{\collapseCVmed}{0.60}
\newcommand{\collapseQbest}{85}
\newcommand{\collapseCVbest}{0.29}
\newcommand{\collapseCV}{0.042}
\newcommand{\collapseD}{0.0628}
\newcommand{\collapseDsd}{0.0029}
\newcommand{\collapseIoU}{0.882}
\newcommand{\collapseCVten}{0.10}
\newcommand{\collapseCVqtr}{0.28}
\newcommand{\tailGainTop}{2.46}
\newcommand{\tailSeTop}{9.11}
\newcommand{\tailNetTop}{0.27}
\newcommand{\tailNetMid}{0.42}
\newcommand{\stNtop}{337}
\newcommand{\rGenAll}{0.095}
\newcommand{\rDsecTrain}{1.31}
\newcommand{\ddRatioMed}{13.9}
\newcommand{\ddBoxesTrain}{290\,163}
\newcommand{\qGenAll}{0.70}
\newcommand{\qDsecTrain}{5.59}
\newcommand{\qDsecRatio}{7.97}
\newcommand{\qDsecNet}{25.7}

%% The recurrent influence tail (E59), from experiments/e42_recurrent_support/tail.json.
\newcommand{\tailExp}{0.645}
\newcommand{\tailExpSE}{0.005}
\newcommand{\tailExpLo}{0.636}
\newcommand{\tailExpHi}{0.655}
\newcommand{\tailRatio}{1.864}
\newcommand{\tailRatioLo}{1.851}
\newcommand{\tailRatioHi}{1.876}
\newcommand{\tailOutside}{77.3}
\newcommand{\tailRmsPower}{2.5}
\newcommand{\tailRmsExp}{23.8}
\newcommand{\tailRmsRatio}{9.4}

%% Per-model cells of the supplement's chunk-position table (E58, all labelled boxes).
%% Three decimals: the smallest difference the table reports is 0.014 pt, so at one
%% decimal the short and full columns would not reproduce the raw column they are
%% differenced into. At three they close against the printed 2-decimal raw values.
\newcommand{\chunkShortSmall}{31.578}
\newcommand{\chunkFullSmall}{36.410}
\newcommand{\chunkRelSmall}{35.6}
\newcommand{\chunkShortT}{38.180}
\newcommand{\chunkFullT}{38.194}
\newcommand{\chunkShortS}{38.371}
\newcommand{\chunkFullS}{39.167}
\newcommand{\chunkShortB}{41.561}
\newcommand{\chunkFullB}{42.126}
\newcommand{\chunkDidT}{0.67}
\newcommand{\chunkSeT}{0.43}
\newcommand{\chunkDidS}{-0.51}
\newcommand{\chunkSeS}{0.48}
\newcommand{\chunkDidB}{-0.16}
\newcommand{\chunkSeB}{0.55}
\newcommand{\aFitShort}{0.628}
% E58e chunk-position curve (src/e58e_curve.py, src/e58f_fig.py)
\newcommand{\chunkRiseBase}{5.78}
\newcommand{\chunkRiseSmall}{4.83}
\newcommand{\chunkRiseRvt}{0.80}
% E58g position-permutation null (src/e58g_perm.py)
\newcommand{\chunkPermSd}{0.55}
\newcommand{\chunkPermZBase}{9.8}
\newcommand{\chunkPermZSmall}{7.6}
\newcommand{\chunkPermPlaceboZ}{1.52}
\newcommand{\chunkPermP}{0.005}
\newcommand{\chunkPermB}{200}
% E62 association-rule sweep (src/e62_association.py)
\newcommand{\assocRules}{14}
\newcommand{\assocLo}{-8.96}
\newcommand{\assocHi}{-0.90}
\newcommand{\assocSpread}{8.06}
\newcommand{\assocSE}{7.12}
\newcommand{\assocGate}{-2.96}
\newcommand{\assocGateSE}{7.12}
\newcommand{\assocFree}{-5.58}
\newcommand{\assocFreeSE}{7.86}
\newcommand{\assocFreeN}{22\,895}
\newcommand{\assocGateN}{22\,518}
%% the rebuild's own gate against the reported table: |(-2.963)-(-2.40)|/7.123
\newcommand{\assocGateVsTable}{0.08}
\newcommand{\assocDelta}{0.37}
\newcommand{\assocMinLo}{12.01}
\newcommand{\assocMinHi}{16.35}
% E63 checkpoint-port parity (src/e63_port_parity.py)
\newcommand{\portEvalGap}{0.101}
\newcommand{\portParamDev}{0.051}
\newcommand{\portSeqs}{429}
\newcommand{\portApT}{38.95}
\newcommand{\portApOwnT}{38.90}
\newcommand{\portApFifty}{69.7}
% E64 sequence-clustered bootstrap of the newest-window profile (src/e64_seqboot.py)
\newcommand{\occSeqs}{12}
\newcommand{\occSeqCiLo}{-24.07}
\newcommand{\occSeqCiHi}{-23.56}
\newcommand{\occSeqSE}{0.13}
\newcommand{\occSeqWidest}{0.51}
\newcommand{\occSeqLoo}{0.08}
\newcommand{\occPerSeqLo}{-24.63}
\newcommand{\occPerSeqHi}{-23.20}

%% E60, the same-frame chunk-boundary intervention: every chunk boundary moved
%% \ssmPairedShift{} windows later, so the frames the release gave one to four windows of
%% history are re-scored with seventeen to twenty. Within model, within frame, identical
%% ground truth and weights; no control model and no parallel-trends assumption.
%% Cluster bootstrap over the same 406 sequences E56 and E58 use.
%% Both released SSM checkpoints are shifted. Suffixes follow \chunkDid*: Small is
%% s5vit-small, Base is s5vit-base, the checkpoint the differenced headline is quoted on.
\newcommand{\ssmPairedN}{3\,574}
\newcommand{\ssmPairedShift}{5}
\newcommand{\ssmPairedMapVelBase}{5.23}
\newcommand{\ssmPairedMapVelBaseSE}{0.49}
\newcommand{\ssmPairedMapVelSmall}{4.81}
\newcommand{\ssmPairedMapVelSmallSE}{0.52}
\newcommand{\ssmPairedMapAllBase}{4.20}
\newcommand{\ssmPairedMapAllSmall}{4.10}
\newcommand{\ssmPairedConfBase}{0.0579}
\newcommand{\ssmPairedConfBaseSE}{0.0048}
\newcommand{\ssmPairedConfSmall}{0.0509}
\newcommand{\ssmPairedConfSmallSE}{0.0035}
\newcommand{\ssmPairedDpfBase}{3.49}
\newcommand{\ssmPairedDpfSmall}{5.09}

%% Remaining E60 rows, for the supplement's paired table. Both channels move: the boxes
%% that match a label are scored higher, and far fewer boxes are emitted in total.
\newcommand{\ssmPairedB}{300}
\newcommand{\ssmPairedMapAllBaseSE}{0.36}
\newcommand{\ssmPairedMapAllSmallSE}{0.32}
\newcommand{\ssmPairedConfTpBase}{0.0409}
\newcommand{\ssmPairedConfTpBaseSE}{0.0040}
\newcommand{\ssmPairedConfTpSmall}{0.0519}
\newcommand{\ssmPairedConfTpSmallSE}{0.0043}
\newcommand{\ssmPairedDpfBaseSE}{0.25}
\newcommand{\ssmPairedDpfSmallSE}{0.26}
\newcommand{\ssmPairedMapVelBaseZ}{10.7}
\newcommand{\ssmPairedMapVelSmallZ}{9.3}
\newcommand{\ssmPairedConfBaseZ}{12.0}
\newcommand{\ssmPairedConfSmallZ}{14.6}
%% The qualitative pair in Fig. 2, and the population it was chosen from. The four
%% count macros record what the panels themselves draw, so the drawn text and this
%% file are checked against one source. The frame is the
%% clearest single instance, so the per-frame averages are quoted beside it in the caption.
\newcommand{\pairedFigLabels}{4}
\newcommand{\pairedFigMatchShort}{0}
\newcommand{\pairedFigMatchFull}{3}
\newcommand{\pairedFigUnmatchShort}{2}
\newcommand{\pairedFigUnmatchFull}{0}
\newcommand{\pairedFigShow}{0.3}
\newcommand{\pairedFigMeanMatchShort}{1.40}
\newcommand{\pairedFigMeanMatchFull}{1.53}
\newcommand{\pairedFigMoreMatched}{13.1}
\newcommand{\pairedFigFewerMatched}{2.7}

\definecolor{cvprblue}{rgb}{0.21,0.49,0.74}
\usepackage[pagebackref,breaklinks,colorlinks,allcolors=cvprblue]{hyperref}
\def\paperID{*****}
\def\confName{CVPR}
\def\confYear{2027}

\newcommand{\micro}{\ensuremath{\mu}}
\newcommand{\eff}{\ensuremath{\hat{\tau}}}

% Figure assets are generated by src/make_figs.py, fall back to a framed box so
% the source compiles before they are rendered.
\newcommand{\figasset}[2]{%
  \IfFileExists{#1}{\includegraphics[width=#2]{#1}}%
  {\fbox{\begin{minipage}[c][3.4cm][c]{\dimexpr#2-2\fboxsep-2\fboxrule\relax}%
   \centering\small\texttt{#1}\\[2pt]\textit{asset not yet rendered}%
   \end{minipage}}}}

\title{Temporal-Support Identifiability\\
in Event-Detection Evaluation Pipelines}
\author{Anonymous CVPR submission}

\begin{document}
\maketitle

\begin{abstract}
An event-detection benchmark defines a sample by one nominal timestamp and scores a prediction
against the label carried at that instant. A recurrent detector forms that prediction from a
history much wider than the instant names, and the released protocol records no variable for
it. A nominal timestamp may therefore fail to identify the temporal evidence that a reported
score represents. We name that protocol property temporal-support identifiability and examine
it with a controlled intervention.

In the released SSM-ViT evaluation pipeline the recurrent history available to a prediction is
a deterministic function of chunk position and ranges from one to twenty-one windows inside a
single pooled score. Moving only the chunk boundary---same frames, same weights, same labels,
same evaluator, and the evaluated timestamp unchanged---changes S5-B by \ssmPairedMapAllBase{}
mAP over all labeled boxes and \ssmPairedMapVelBase\,$\pm$\,\ssmPairedMapVelBaseSE{} mAP on the
velocity-evaluable subset. The evidence behind the score moved while every quantity the
protocol records stayed fixed.

Two audits delimit that result on either side of the pipeline. On the predictor side, the
newest RVT input window of the ported released checkpoint carries an ablation-sensitivity
centroid \rvtCentroidMagMeas\,ms before the label instant, alongside a regression-based
temporal coefficient of \ladderFullTau\,$\pm$\,\ladderFullTauSE\,ms that we read as an
association measured on nominal-time-matched detections and not as an estimator of recurrent
support. On the label side, within-exposure event timing carries an illumination-locked
component, so label-acquisition support is not attributable to object motion alone. A benchmark
timestamp is thus not by itself an identifiable description of the evidence supporting a
recurrent prediction, and a score reported under one does not license a timing comparison.
\end{abstract}

\section{Introduction}

An event-detection benchmark defines a sample by a nominal timestamp. A prediction $\hat{s}$
is scored against the ground truth $s^{*}$ carried at a query time $t$ through a function
$\rho(\hat{s},s^{*})$ of position, shape and class. Beyond selecting the label at $t$, that
function contains no variable representing the temporal evidence behind either argument, so a
temporal discrepancy enters the score as displacement.

A recurrent detector does not read only the instant the benchmark names. It computes $\hat{s}$
from an input window and from state carried across earlier windows, and the label it is
compared against comes from an exposure of finite width, an annotator or tracker working on the
exposed image, and a geometric transfer into another sensor's frame. The evidence on both sides
is an interval, while the protocol records an instant.

The benchmark timestamp may therefore fail to identify the temporal evidence that a reported
score represents. Three terms carry this throughout. The \emph{nominal timestamp} is the instant
at which the protocol pairs a prediction with a label. The \emph{effective temporal support} is
the past that can influence the prediction, including state carried by recurrence. The
\emph{label-acquisition support} is the interval over which the label's evidence was acquired.
A nominal timestamp determines neither of the other two uniquely, and a temporal descriptor may
be reconstructable from released artifacts yet remain unreported or heterogeneous within one
pooled score. We formulate the resulting protocol property as temporal-support identifiability:
whether the released protocol and artifacts uniquely determine the temporal quantities a
reported score represents.

We examine it with a paired intervention on the released SSM-ViT evaluation pipeline. In that
implementation the nominally carried state is inert across evaluation chunks, so the recurrent
history available to a prediction is a deterministic function of chunk position and ranges from
one to twenty-one windows inside one pooled score. Moving every chunk boundary
\ssmPairedShift{} windows later re-scores the same frames under a longer history while frames,
weights, ground truth, evaluator and the evaluated timestamp are held fixed
(\ssmPairedN{} paired boxes, cluster bootstrap over \outSeqsScanned{} sequences).

Effective temporal support alone changes S5-B by \ssmPairedMapAllBase{} mAP over all labeled
boxes and \ssmPairedMapVelBase\,$\pm$\,\ssmPairedMapVelBaseSE{} mAP on the velocity-evaluable
subset (Sec.~\ref{sec:chunk}); the second released checkpoint moves by
\ssmPairedMapAllSmall{} and \ssmPairedMapVelSmall{} points. Every quantity the protocol records
stayed fixed, and the score moved. A benchmark timestamp is thus not by itself an identifiable
description of the evidence supporting a recurrent prediction.

Work on event-detection timing takes that index as given and improves what is measured against
it: sensor latency is calibrated away~\cite{yang2024latency, liu2024nighttime}, compute latency
is folded into average precision~\cite{li2020streaming,yang2022streamyolo,chu2026stare}, and the
temporal window is redesigned~\cite{gehrig2023rvt,zubic2024ssm}. Each changes what the predictor
does at a fixed benchmark index, whereas the object measured here is that index itself and the
support a score reported under it pools. The two datasets below instantiate this protocol
condition; how often it holds elsewhere is not measured.

\textbf{Supporting audit: predictor-side support.} The event branch of the detector examined
here~\cite{gehrig2023rvt} consumes a window of \rvtWindow\,ms in \rvtBins{} bins whose end sits
on the label instant (Sec.~\ref{sec:window}), so under uniform weighting its evidence is
centered \rvtCentroidMag\,ms earlier. The trained network's bin sensitivities are not uniform,
and ablating each bin in turn on real validation samples puts the ablation-sensitivity centroid
at \rvtCentroidMeas\,ms (Sec.~\ref{sec:influence}). Over a 1000\,ms horizon the newest window
accounts for 22.7\% of the summed window-ablation sensitivity and the remaining 77.3\% lies in
preceding windows. The regression-based temporal coefficient, estimated on nominal-time-matched
detections under Eq.~2 with a vector regression, a placebo term and a sign-flip control on the
velocity estimator, is \ladderFullTau\,$\pm$\,\ladderFullTauSE\,ms and sits \outInterval\,ms
from that centroid under the reported specification. It is compatible with zero and is reported
as a regression association between nominal-time-matched detections, not as an estimate of the
network's recurrent support; the two are distinct operational quantities the nominal timestamp
does not encode (Secs.~\ref{sec:window}--\ref{sec:extrapolation}).

\textbf{Supporting audit: label-acquisition support.} Across the \dsecSeqCount{} DSEC training
sequences whose exposure files resolve, the published exposure width varies by a factor of
\dsecExpRatio{} against a frame period fixed at \dsecFramePeriod\,ms, with \dsecPinnedFrac\,\%
of frames pinned at an auto-exposure ceiling in \dsecCeilingSeqs{} sequences
(Fig.~\ref{fig:exposure}). Those sequences carry the harmonic series of a \flickerLineHz\,Hz
mains flicker, and between \sixLockedLoPct\,\% and \sixLockedHiPct\,\% of labeled boxes exceed
the ninety-ninth percentile of the off-frequency local excess measured on the same events.
Within-exposure event timing therefore contains illumination-locked structure and is not
attributable to object timing without separating it (Secs.~\ref{sec:flicker}--\ref{sec:cost}).

Identifiability under the metric is then read in the metric's own units. On the
velocity-evaluable subset, displacing ground-truth centers across the measured label-to-centroid
interval changes mAP by \mapCost{} points, and at that displacement the five released
checkpoints are not reordered. Separately, over all labeled boxes the nominal validation-split
ordering is reproduced in only \rankHoldBox\,\% of sequence-bootstrap replicates. These results
place the observed displacement effect at the benchmark's sampling-resolution floor.
Table~\ref{tab:strata} relates this resolution to the benchmark's speed composition (Supplement
Secs.~11--12).

\section{Related work}
\label{sec:related}

\paragraph{Sensor and runtime latency.}
Prior work calibrates sensor latency
\cite{yang2024latency,liu2024nighttime}, incorporates runtime latency into
average precision \cite{li2020streaming,yang2022streamyolo,chu2026stare}, or
proposes temporal tolerance \cite{softed2024}. These methods address the time
from an event to an output and take the benchmark's temporal index as given.
We instead measure that index, testing whether a released benchmark
specification uniquely determines the temporal quantities behind its score.

\paragraph{Temporal windows and benchmark identifiability.}
RVT and SSM-ViT redesign or parameterize temporal windows
\cite{gehrig2023rvt,zubic2024ssm}, while adaptive-window methods compare
alternative partitions \cite{sui2026astw}. Event--RGB frequency mismatch and
exposure duration have also been noted \cite{faod2024,sun2023refid,kim2022unknownexposure}.
The closest accepted work above calibrates sensor or runtime latency or redesigns temporal
windows, each holding the benchmark's temporal index fixed. We make that index the measured
quantity, reading the released per-frame windows and exposures and separating predictor
support, output-time interpretation, and label-acquisition support.

\paragraph{Benchmark reliability and evaluation analysis.}
Prior work studies detector calibration, stratified model evaluation, label-induced benchmark
instability, and event-vision evaluation protocols
\cite{oncalibration2024,stratification2024,northcutt2021labelerrors,ercan2023evreal}.
These analyses examine uncertainty or error in reported performance, whereas our focus is the
temporal support and acquisition semantics that a nominal detection timestamp leaves implicit.

\paragraph{Label time and attribution.}
Waymo and annotation studies estimate per-object capture or annotation offsets
\cite{waymolabelproto,sun2020waymo,nunez2026annotations}. Event datasets expose
per-box time fields, but their degrees of freedom are not necessarily specified.
LET-3D-AP and TIDE/DETAD/HOTA decompose or bucket spatial error while taking the
label clock as given \cite{hung2022let,bolya2020tide,hoiem2012diagnosing,alwassel2018detad,luiten2021hota}.

\paragraph{Event benchmarks and exposure.}
Gen1, 1\,Mpx, DSEC-Det and Ev-3DOD use asynchronous measurements, rectified
frames, or interpolated labels \cite{detournemire2020gen1,perot2020megapixel,
gehrig2024lowlatency,cho2025ev3dod}. LEOD establishes causal input windows
\cite{wu2024leod}. The present work reads the released window and exposure
artifacts as protocol evidence and tests whether the resulting temporal
quantities are uniquely specified.

\iffalse
\paragraph{Recovering a time offset from image-plane displacement.} Qin and
Shen~\cite{qin2018temporal} model a global camera--IMU time offset by shifting
feature observations with image-plane velocity and jointly optimizing the offset
with visual--inertial states and feature locations. Here the offset lies between
a predictor's clock and
a \emph{benchmark's} label clock. Streaming perception~\cite{li2020streaming}
folds wall-clock compute latency into average precision and re-anchors any
detector from box coordinates alone. Later work adapts to the measured runtime
delay~\cite{yang2022streamyolo}, a continuous-stream framework scores the latest
available output against high-rate ground truth and reports reordered
rankings~\cite{chu2026stare}, and temporal tolerance has been argued for
time-series event detection~\cite{softed2024}. Each of those latencies is the
time from an event to an output and falls as the device gets faster. Sensor
latency, the illumination-dependent delay between a change and the timestamp the
sensor assigns it, is measured at
microseconds~\cite{yang2024latency}, and nighttime event imaging corrects an
illumination-dependent trailing by learning a per-event timestamp
calibration~\cite{liu2024nighttime}. Both target the sensor, while here the
released exposure metadata and event stream bound what a within-box temporal
statistic computed inside one published exposure can be attributed to. Inside
the citation graph of LET-3D-AP~\cite{hung2022let}, a latency-aware average
precision scores 3D detection at the instant a prediction becomes available on
a named device~\cite{rethink3d2025}. That work addresses compute latency,
whereas the interval measured in Sec.~\ref{sec:extrapolation} is fixed by a data representation. The
event--RGB frequency mismatch is named in~\cite{faod2024}, which scores with
unmodified average precision.

Mao et al.~\cite{mao2019delay} showed that video-detection AP can remain stable
despite temporal detection delay and introduced Average Delay to expose that
effect. This work instead asks which acquisition and label times a nominal benchmark
timestamp represents, using released detector and exposure artifacts.

\paragraph{Decomposing localization error along a privileged axis.}
LET-3D-AP~\cite{hung2022let} decomposes camera-only 3D localization error along
the line of sight and defines an average precision tolerant of the longitudinal
part, and reports a camera-to-label synchronization gap as a per-dataset interval
that depends on horizontal object position. Sec.~\ref{sec:outputtime} keeps the
idea of a privileged axis but replaces the sensor's line of sight with the
object's own velocity vector, and estimates a time rather than defining a
tolerant score.
TIDE~\cite{bolya2020tide,hoiem2012diagnosing} and DETAD~\cite{alwassel2018detad}
bucket detection error atemporally, so a temporal offset falls inside their
localization bucket, and HOTA~\cite{luiten2021hota} takes the label clock as
exact.

\paragraph{Per-object time in labels.} Waymo's \texttt{camera\_synced\_box} solves
$t_{\mathrm{capture}}$ per object and stores the resulting box shift rather than
the time~\cite{waymolabelproto,sun2020waymo}. Miro et
al.~\cite{nunez2026annotations} measure within-frame per-object annotation
offsets on Argoverse\,2 and TruckScenes and define a dispersion of the
\emph{spatial} error. Each computes a per-object time in order to remove it. The
event-vision label formats carry the field: the Prophesee dtype used by Gen1 and
the 1\,Mpx dataset, and the DSEC-Det format, both have a per-box \texttt{t}. In the releases examined
here, that field's per-object degree of freedom is not specified. Across the \ddSeqs{} released DSEC-Det
sequences, \ddBoxes{} boxes take \ddDistinctT{} distinct values of \texttt{t},
\ddBoxesPerT{} boxes per value, and those values are the \dsecFramePeriod\,ms
frame clock.

\paragraph{Event benchmarks and their exposures.}
Gen1~\cite{detournemire2020gen1} labels are drawn on gray-level images assembled
from asynchronous per-pixel measurements. On the 1\,Mpx
dataset~\cite{perot2020megapixel} the event count inside a labeled box at its own
timestamp is often too small to support detection~\cite{kugele2023howmany}.
DSEC-Det~\cite{gehrig2024lowlatency} labels by tracking~\cite{pang2021qdtrack} on
rectified frames and warps into the event view, and
Ev-3DOD~\cite{cho2025ev3dod} interpolates lower-rate annotations into high-rate
3D labels. That a frame's exposure is long relative to event timing has been
stated: Event-based deblurring~\cite{sun2023refid} argues that accumulating events from a single
frame timestamp loses information, deblurring work names auto-exposure as why the
width is unknown~\cite{kim2022unknownexposure}, and RENet~\cite{zhou2022renet}
uses a DSEC exposure interval as its event window while describing it as short.
The cited works do not read the released per-frame widths or measure what is
inside one. State space backbones~\cite{zubic2024ssm} substantially reduce sensitivity to
inference-time window duration, and
adaptive spatial-temporal windows~\cite{sui2026astw} retrain a detector under
each of nine partitioning strategies and report the score of each. Both treat the
window as a design variable to choose. The released window configuration is
treated as a measured property of the released detector and evaluation pipeline.
That the released window is causal is already in print:
LEOD~\cite{wu2024leod} states that event detectors take only events triggered
before the query time, and builds a time-flip augmentation on the consequence
that forward inference favors receding objects. It does not measure where inside
that window the evidence sits. Protocol
work supplies the standards these measurements are held
to~\cite{oncalibration2024,stratification2024,northcutt2021labelerrors,ercan2023evreal}:
a repaired protocol can reverse a field's conclusion, and a reported score is
an estimate with a variance. In this benchmark, the streaming intervention changes the
reported score even though the model and evaluated frames are held fixed.

\fi
\section{Effective temporal support under a released streaming evaluation}
\label{sec:chunk}

The state the SSM-ViT release carries between evaluation chunks does not affect any of its
outputs: it enters as a multiplicative factor on the first coefficient of an associative scan
on which the outputs do not depend. Under the released streaming evaluation this bounds the
recurrent history available to a detection to the current \ssmEvalChunk{}-window chunk, so that
history ranges from one window to \ssmEvalChunk{} with the detection's position in the chunk,
and the reported number averages over all \ssmEvalChunk{} positions. Chunk position is a
deterministic function of the release's own index files and is recorded nowhere in the
evaluation protocol.

\paragraph{The same-frame intervention.}
The intervention moves every chunk boundary \ssmPairedShift{} windows later and changes nothing
else: the same frames, the same weights, the same labels, the same evaluator, and the evaluated
timestamp of each frame unchanged. \ssmPairedN{} labeled boxes are scored twice under this
shift, once with one to four windows of recurrent history and once with seventeen to twenty,
and are compared paired within frame with a cluster bootstrap over the \outSeqsScanned{}
sequences. The long-history arm gains
\ssmPairedMapVelBase$\pm$\ssmPairedMapVelBaseSE{} mAP points for S5-B and
\ssmPairedMapVelSmall$\pm$\ssmPairedMapVelSmallSE{} for S5-S on the velocity-evaluable subset,
and over every labeled box the S5-B gain is
\ssmPairedMapAllBase$\pm$\ssmPairedMapAllBaseSE{} points. Mean detection confidence, which uses
no labels, moves with it by \ssmPairedConfBase$\pm$\ssmPairedConfBaseSE{} and
\ssmPairedConfSmall$\pm$\ssmPairedConfSmallSE{}. Fig.~\ref{fig:chunkpos} (bottom) draws one such
pair. No quantity the protocol records distinguishes the two arms.

\paragraph{The observational contrast across chunk positions.}
Reconstructing chunk position over \mapFrames{} labeled frames of the released evaluation and
scoring every labeled box, S5-B reaches \chunkShortBase{} mAP on the frames given one to four
windows and \chunkFullBase{} on those given seventeen to \ssmEvalChunk{}
(Fig.~\ref{fig:chunkpos}, top). This contrast compares different frames, so it is differenced
against the three RVT checkpoints, evaluated on the same frames with a state that crosses chunk
boundaries; that controls for frame-position effects shared with them and leaves
\chunkDidBase$\pm$\chunkDidBaseSE{} points for S5-B ($z=\chunkDidBaseZ$) and
\chunkDidSmall$\pm$\chunkDidSmallSE{} for S5-S ($z=\chunkDidSmallZ$), against placebo
differences among the RVT checkpoints of at most \chunkPlaceboMax{} points
($z=\chunkPlaceboZ$). Permuting chunk position within each sequence puts both SSM values at
least \chunkPermZSmall{} standard deviations outside their own null and every placebo inside
it. The two readings agree in sign and magnitude, and the paired intervention carries the claim
because it holds the frames fixed (Supplement Sec.~13).

\paragraph{What the pooled score reports.}
One released number therefore pools detections whose recurrent history differs by a factor of
\ssmEvalChunk{}. Restricting evaluation to the frames receiving 17--21 windows yields an S5-B
score \chunkGapBase{} points above the pooled released-evaluation score. Two systems compared
under this protocol at the same nominal timestamps need not have been given the same temporal
evidence, and the protocol records no variable by which a reader could tell.

\begin{figure}[t]
\centering
\figasset{figs/chunkpos.pdf}{\columnwidth}
\\[1.5pt]
\figasset{figs/paired_frame.pdf}{\columnwidth}
\caption{\textbf{Top:} benchmark score against the recurrent history the released streaming
evaluation made available, at each of the \ssmEvalChunk{} chunk positions, over every
labeled Gen1 box, no displacement applied. Shaded: the one-to-four-window block the
release scores together with the rest. Between that block and the frames given seventeen
windows or more, the two SSM checkpoints rise by \chunkRiseBase{} and \chunkRiseSmall{}
points and the three RVT checkpoints, scored on the same frames with state that crosses
chunk boundaries, by at most \chunkRiseRvt{}.
\textbf{Bottom:} one frame of the same-frame intervention, S5-B. Both panels share the
frame, the labels (dashed) and the weights, and only the chunk boundary moves. Solid: a
detection matched to a label of its class at IoU $\ge 0.5$. Dotted: unmatched. Drawn above
confidence \pairedFigShow{} and scored at every confidence. The example is illustrative only.
All paired quantitative estimates use all \ssmPairedN{} frames. Supplement Sec.~13.}
\label{fig:chunkpos}
\end{figure}

\section{Supporting audit: predictor-side temporal descriptors}
\label{sec:predictor}

\subsection{Identifiability and the timing estimands}

A benchmark releases one tuple per evaluated instant: a nominal timestamp, the
labels asserted at it, the predictions scored against them, and the score. Three
temporal quantities are latent in that tuple rather than recorded in it: the temporal
support of the predictor's input and carried state, the effective state time its output
refers to, and the acquisition support of the labels. Call such a quantity
\emph{protocol-identifiable} when the released protocol and artifacts determine it uniquely.
Under this definition, input and exposure intervals are protocol-identifiable, whereas
effective predictor timing, recurrent context represented by a pooled score, and
within-exposure attribution are not.

Let a predictor $P$ emit a state at query time $t$. A signed effective-weight
model uses bin centers $c_k$ and weights $a_k$ with $\sum_k a_k=1$, giving the
effective timestamp offset $\tau_P=\sum_k a_kc_k$ under constant velocity. The
empirical ablation profile instead assigns non-negative sensitivities $s_k\geq0$
to bins and has centroid $c_s=\sum_k s_kc_k/\sum_k s_k$. These are distinct
quantities: zero-ablation sensitivity depends on both the predictor and the
input, and is not in general a decomposition of temporal weights. Sections~\ref{sec:window}--\ref{sec:influence} measure this
operational input-sensitivity profile over the newest window. Section~\ref{sec:outputtime}
estimates the regression coefficient $\tau$ under Eq.~2. Interpreting $-\tau$ as an
effective state-time offset requires the absence of residual velocity-proportional
spatial bias.
Section~\ref{sec:extrapolation} compares $-c_s$ with the regression coefficient
$\tau$, equivalently $c_s$ with the regression-implied timestamp $\tau_P=-\tau$.

\subsection{The input window and its indexing instant}
\label{sec:window}

Two released files fix the input side. The dataset configuration
\texttt{config/dataset/gen1.yaml} of the detector's own
repository~\cite{gehrig2023rvt} names the representation
\texttt{stacked\_histogram\_dt=50\_nbins=10}, a window of \rvtWindow\,ms divided
into \rvtBins{} bins of \rvtBinWidth\,ms. Line \rvtSourceLine{} of
\texttt{scripts/genx/preprocess\_dataset.py} constructs the array
\texttt{ev\_repr\_timestamps\_us\_end} by counting backwards from a label
timestamp in steps of the window duration, and extends it between consecutive
label timestamps. Each representation is therefore indexed by the end of its
window, and those ends are placed on label times. A window attached to a label at
time $t$ holds the events of $[t-\rvtWindow\,\mathrm{ms},\,t]$, whose uniform
weighting centroid falls \rvtCentroidMag\,ms before the instant $t$ at which the
label asserts a state. We read both statements from the released source, and
the \rvtCentroidMag\,ms figure is a reference point rather than a property of
the trained network.

\begin{figure}[t]
\centering
\figasset{figs/fig_predictor.pdf}{0.95\columnwidth}
\caption{Input sensitivity and the regression-based temporal coefficient summarize the same
nominal timestamp without being equivalent, and are not competing estimates of a common latent
timestamp; their point estimates are \outInterval\,ms apart. Panel
(b) shows that average precision on the velocity-evaluable subset is nearly
invariant over the measured \rvtCentroidMagMeas\,ms
label-to-centroid displacement. The \outInterval\,ms separation is not an estimate of
architectural latency.
(a) Window-ablation influence of each of the \rvtBins{} bins of the released
\texttt{rvt-t} input window, over \rvtInfluenceSamples{} Gen1 validation samples with the
recurrent state warmed \rvtWarmSteps{} steps, with sample-level standard errors. Error bars are
descriptive standard errors across the sampled instances and are not used for inference.
Uncertainty of the ablation-sensitivity centroid is evaluated by sequence bootstrap in
Supplement Sec.~2. Dashed:
the mean over bins. Dotted: the label instant. Solid: the ablation-sensitivity centroid.
Dash-dot, with its standard error shaded: the regression-implied timestamp under Eq.~2,
$\tau_P=-\tau$, of Sec.~\ref{sec:outputtime}.
Gray: the four other checkpoints, each rescaled to its own mean. (b) Average precision
against the temporal center displacement $\delta$ applied to the ground-truth boxes, by
object speed in px/s, each curve relative to its own value at the label instant, with its
maximum marked. The shaded band runs from the centroid to the label instant, spanning
\outPxNinety\,px at the ninetieth percentile of label speed against an annotation
dispersion of \labelSigmaGenAnn\,px (Supplement Sec.~3). The mAP span is \stSpanAll{}
points over all moving boxes and \mapSpanHigh{} points in the
\mapCutMid--\mapCutHi\,px/s stratum, where the centroid displacement costs
\stCostHigh{} points.}
\label{fig:predictor}
\end{figure}

\subsection{The measured influence profile}
\label{sec:influence}

The newest-window ablation-sensitivity centroid is computed from the measured
bin-ablation sensitivities of a released checkpoint. We take \rvtInfluenceSamples{} samples of the released Gen1
validation split with the released \texttt{rvt-t} checkpoint, warm the recurrent
state for \rvtWarmSteps{} steps, and ablate each bin in turn. For sample $i$, let
$Y_i$ denote the emitted detection tensor from the unablated input and
$Y_i^{(-k)}$ the tensor obtained after zeroing temporal bin $k$ while keeping the
incoming recurrent state identical. We define the ablation sensitivity, or influence, of bin $k$ as
\begin{equation}
\begin{aligned}
s_k&=\frac{1}{N}\sum_{i=1}^{N}
\frac{\left\|Y_i^{(-k)}-Y_i\right\|_2}{\left\|Y_i\right\|_2},\quad N=480,\\
c_s&=\frac{\sum_k s_k c_k}{\sum_k s_k},
\end{aligned}
\end{equation}
where $c_k$ is the center time of bin $k$ relative to the label instant. The
ablated forward pass receives the same incoming recurrent state as the
corresponding reference pass, and its returned state is discarded. The profile is
a property of the current window and not of a perturbed history.

The profile is balanced between the two halves of the window (Fig.~\ref{fig:predictor}a)
and yields an ablation-sensitivity centroid of \rvtCentroidMeas\,ms, compared
with \rvtCentroidUnif\,ms for uniform weighting. The released configuration
therefore places the measured centroid near the window reference without implying
uniform network sensitivity.
Across zero, mean, cross-sample replacement, and gradient instruments, the centroid spans
\occMean{} to \occGrad{} ms (Supplement Sec.~2).
The reported centroid applies to the deterministic 12-sequence subset described in the supplement.
This is the current window's window-ablation sensitivity profile with the recurrent
state held fixed, not a decomposition of the full temporal support.

\paragraph{Recurrent support beyond the newest window.}
Over a 1000\,ms horizon, the newest window accounts for 22.7\% of the summed
window-ablation sensitivity, while the remaining 77.3\% lies in preceding windows. The
horizon centroid shifts from
\histCentroidTen\,ms at \histHorizonA\,ms to \histCentroidTwenty\,ms at \histHorizonB\,ms,
so the measured range does not determine a converged full-history centroid. Supplement Sec.~14
characterizes the tail and its fill controls.

\paragraph{Across the released checkpoints and a second architecture.}
RVT releases three Gen1 checkpoints differing only in embedding width, attention
head dimension and FPN depth. SSM-ViT~\cite{zubic2024ssm} releases two more in
which each stage's ConvLSTM is replaced by an S5 state space layer, at \ssmParamsS{}
and \ssmParamsB{} M parameters against RVT's \capParamsS{} and \capParamsB{}.
It distributes RVT's own preprocessed Gen1 and inherits its window construction. On the same sequences, and
with each driven as its own release specifies, all \archCkpts{} checkpoints put the
newest window's ablation-sensitivity centroid between \archHi{} and \archLo\,ms, none further than
\archUnifMax\,ms, or \archUnifPct\,\% of one bin, from the uniform-weight value.

The support within the newest window
is not fixed across capacities: over \capSamples{} samples of the
three RVT capacities, the newest window's share of a \capHorizon\,ms horizon
falls from \capShareT\,\% to \capShareB\,\% and the sensitivity centroid moves
from \capSuppT{} to \capSuppB\,ms.

\subsection{Regression-based temporal coefficient}
\label{sec:outputtime}

The influence profile is a statement about the network's input. Even a predictor
with uniform local sensitivity can emit a state describing a later instant,
and that quantity is measured from the same checkpoint's own detections.

A detector reporting the state an object occupied $\tau$ earlier incurs the
physical displacement $-\tau v$, where $v$ is the latent true velocity. We
observe only the finite-differenced label velocity $\tilde v=v+\eta$ and
therefore fit Eq.~\ref{eq:outreg} using $\tilde v$. Thus
$\tau=-\tau_{P}$ on the axis of Sec.~\ref{sec:predictor}. With $\hat{c}$ the
predicted center, $c^{*}$ the label center, $A$ the area of the ground-truth
box, $\hat{u}=\tilde v/\lVert \tilde v\rVert$ and $\mathrm{R}$ a quarter turn,
$b$ is a two-dimensional intercept, fitted globally or as per-sequence
intercepts in the reported specification,

\begin{equation}
\hat{c}-c^{*}=b-\tau\,\tilde v+\rho\,\mathrm{R}\tilde v+\gamma\sqrt{A}\,\hat{u}+\varepsilon
\label{eq:outreg}
\end{equation}

\noindent is fitted by stacking both components of every match into one least
squares problem, with standard errors clustered by sequence. It is the temporal
calibration of~\cite{qin2018temporal} carried to a box center, taken as a slope
rather than as the per-object ratio $-\langle
\hat{c}-c^{*},\hat{u}\rangle/\lVert \tilde v\rVert$, whose variance diverges as its
denominator approaches zero. For comparison, the per-object ratio evaluated on
the same rows has a mean of
$\ratioMean\,\pm\,\ratioMeanSE$\,ms, on the hypothesis evaluated in
Sec.~\ref{sec:extrapolation}, with an interquartile range of \ratioIQR\,ms.
Written this way, Eq.~2 fits the two-dimensional residual directly with $\tilde v$ as a
regressor. Unlike the per-object ratio, it does not first project the residual onto $\hat u$
and divide by $\|\tilde v\|$.
The reported full specification adds four coordinate/size regressors and
per-sequence intercepts to the terms displayed in Eq.~\ref{eq:outreg}.
The coefficient is a temporal offset only if no velocity-proportional spatial
bias remains after these controls. Such a bias is indistinguishable from $\tau$ in Eq.~2.

$\mathrm{R}\tilde v$ is the label velocity rotated a quarter turn. Because a pure
temporal offset acts along $\tilde v$ rather than $\mathrm{R}\tilde v$, $\rho$ serves
as a negative-control coefficient. A nonzero value implies residual spatial
misspecification.
$\sqrt{A}\,\hat{u}$ absorbs a displacement along travel proportional to box side,
which is needed because that residual grows with box size and box size correlates
with label speed at \outSpeedSizeCorr{} on this split.

\paragraph{The finite difference used for the velocity.}
$\tilde v$ is finite differenced from labels. Label noise enters forward and backward
velocity differences with opposite shared-noise signs. The centered difference excludes the
current label center and removes this direct same-frame term. Supplement Secs.~1 and 4 derive
and calibrate this effect, so the reported regression uses the centered estimator.

\begin{table}[t]
\centering\footnotesize
\setlength{\tabcolsep}{3.2pt}\renewcommand{\arraystretch}{0.95}
\begin{tabular}{lrr}
\toprule
Specification & $\tau$ (ms) & placebo $\rho$ (ms) \\
\midrule
\multicolumn{3}{l}{\emph{velocity estimator, no geometric controls}}\\
\quad forward difference   & \sfMinFwd\,$\pm$\,\sfMinFwdSE   & \sfMinPlFwd\,$\pm$\,\sfMinPlFwdSE \\
\quad backward difference  & \sfMinBack\,$\pm$\,\sfMinBackSE & \sfMinPlBack\,$\pm$\,\sfMinPlBackSE \\
\quad centered difference  & \sfMinCen\,$\pm$\,\sfMinCenSE   & \sfMinPlCen\,$\pm$\,\sfMinPlCenSE \\
\midrule
\multicolumn{3}{l}{\emph{controls added, centered difference throughout}}\\
\quad $+$ box side along travel & \ladderMinTau\,$\pm$\,\ladderMinTauSE  & \ladderMinPl\,$\pm$\,\ladderMinPlSE \\
\quad $+$ four geometric terms  & \ladderGeoTau\,$\pm$\,\ladderGeoTauSE  & \ladderGeoPl\,$\pm$\,\ladderGeoPlSE \\
\quad $+$ per-sequence bias     & \ladderFullTau\,$\pm$\,\ladderFullTauSE& \ladderFullPl\,$\pm$\,\ladderFullPlSE \\
\bottomrule
\end{tabular}
\caption{The regression-based temporal coefficient $\tau$ of the released \texttt{rvt-t} checkpoint
over \outMatches{} matched boxes of \outSeqs{} Gen1 validation sequences, and the
placebo $\rho$ beside it, the coefficient on the rotated label velocity, which
no temporal offset can reach. Standard errors are clustered by sequence. The
forward and backward fits above are calibrated in Supplement Sec.~4.}
\label{tab:outtime}
\end{table}

\paragraph{Regression dataset and matching.}
Rows are detections matched to nominal-time ground truth at IoU $\geq 0.3$ with centered
label velocities available, leaving \outMatches{} matches from \outSeqs{} sequences. The
estimate is conditional on this matched subset. Supplement Sec.~5 reports association-rule
sweeps, including overlap-free matching.

Table~\ref{tab:outtime} is a specification ladder, and the specification was
chosen by the placebo rather than by $\tau$. Under Eq.~\ref{eq:outreg} with only
the box-side term, $\tau=\ladderMinTau\,\pm\,\ladderMinTauSE$\,ms, but $\rho$ is
\ladderMinPl\,ms at \ladderMinPlSigmas{} standard errors, which indicates a term
omitted from Eq.~\ref{eq:outreg} that $\mathrm{R}\tilde v$ is a proxy for. A split by
the sign of the horizontal velocity locates it: $\rho$ changes sign across that
split, at \plRight\,ms and \plLeft\,ms, which a rotational term would not do.
Most label motion here is horizontal, so $\mathrm{R}\tilde v$ maps a horizontal speed
onto a vertical residual. The rotated regressor represents a bias that differs
between oncoming and same-direction traffic. Entering
box height and image row along $\hat{y}$ and box width and image column along
$\hat{x}$, and then a constant vector per sequence, returns the placebo to
\ladderFullPl\,ms at \ladderFullPlSigmas{} standard errors. That is the
specification reported, and the criterion that selected it is a coefficient on a
regressor orthogonal to $\tilde v$.

Under the reported specification, $\tau=\ladderFullTau\,\pm\,\ladderFullTauSE$\,ms, with a
sequence-bootstrap 95\% interval of [\outBootLo,\,\outBootHi]\,ms. Across geometric-control
specifications, $\tau$ spans \satLo{} to \satHi{} ms. Because the reported specification was
selected using the orthogonal placebo coefficient, its nominal standard error does not include
specification-selection uncertainty. We therefore treat the point estimate descriptively and use
the full specification sweep to assess robustness.

\subsection{Relating input-sensitivity timing to the regression-based temporal coefficient}
\label{sec:extrapolation}

The comparison combines measurements under different recurrent-state conditions.
The centroid uses an eight-step warm-up with the incoming state fixed during
ablation. The regression uses detections from the validation stream with the
released state carried across frames. The \outInterval\,ms numerical separation is descriptive only: the two quantities are
distinct operational summaries and are not competing estimates of a common latent timestamp.
Under the reported specification and its
annotation-error assumptions, the point-estimate difference is \outInterval\,ms. The regression
coefficient has a standard error of \outIntervalSE\,ms. This is a regression-coefficient
result rather than a latency estimate of this architecture. Expressed on Eq.~\ref{eq:outreg}'s
sign convention, the newest-window centroid corresponds
numerically to $\tau=+\rvtCentroidMagMeas$\,ms. This comparison relates two operational
estimands and does not assume that Eq.~\ref{eq:outreg} must recover the ablation centroid.
Thus, \rvtCentroidMagMeas\,ms is the measured interval between the label
instant and the newest-window ablation-sensitivity centroid, whereas \outInterval\,ms
is the point-estimate difference between that centroid and the reported regression
coefficient. The measured horizons do not determine a converged full-history centroid. The variation across specifications is substantial
relative to the reported standard error and is not captured by that standard error.

The detector is trained against the ground truth at the label time on a window
ending there, so a measured offset near zero is consistent with the training
objective indexed at the label time. The objective
does not enforce the
interval, and the instant the output describes and the centroid of the evidence
used to produce it are different quantities that one benchmark timestamp stands
for. At the ninetieth percentile of label speed on this split, \outSpeedNinety\,px/s, the
\rvtCentroidMagMeas\,ms is \outPxNinety\,px of displacement, comparable to the
annotation dispersion of the labels themselves (\labelSigmaGenAnn\,px), and at the
ninety-ninth it is \outPxNinetynine\,px. That ablation-sensitivity centroid and
the regression-based temporal coefficient are measured from released artifacts alone.


\paragraph{Effect on benchmark mAP.}
This experiment is not a test of whether temporal support matters in general. It asks whether
Gen1's particular speed composition makes the measured RVT interval visible to its current mAP.
We translate each ground-truth box
center by $\delta\tilde v$, keeping its size and class fixed, a
first-order surrogate for temporal displacement under local constant velocity.
The released predictions remain unchanged, since a detection
has no velocity of its own. Over \mapFrames{} frames, \mapGTvelPct\,\% of the \mapGT{}
ground-truth boxes carry a centered velocity. Ground-truth
boxes without a centered velocity are marked ignored during evaluation, and detections
matched to ignored boxes are excluded from TP/FP accounting. The resulting mAP against
$\delta$ is smooth with its maximum at \mapArgmax\,ms
(Fig.~\ref{fig:predictor}b). Along the axis of Eq.~\ref{eq:outreg} that is
$\tau=\mapArgmaxTau$\,ms, and mAP varies by at most \mapCost{} points over the
interval. This descriptive curve does not by itself identify a preferred temporal
value.

The full $-60$ to $+60$\,ms sweep moves mAP by \mapSweepSpan{} points. On the
$-50$ to $+30$\,ms grid summarized in Table~\ref{tab:strata}, the all-moving
span is \stSpanAll{} points. Displacing the ground truth to the window's centroid
costs \mapCostAbs{} absolute mAP (\mapCost{} percentage points,
\mapCostPct\,\% relative to \mapZero{}). Average precision on the velocity-evaluable
subset is nearly invariant over the measured label-to-centroid displacement, and at
the common $-\rvtCentroidMagMeas$\,ms displacement the ordering of all
\rankCkpts{} released checkpoints is unchanged on the validation split. Scoring each checkpoint at its own
preferred displacement likewise leaves the validation-split ordering unchanged in the reported
strata. Both statements are bounded by the stability of the
ordering under sequence resampling at this scale. The smallest timing difference that would reorder an adjacent
pair is \dstarAll\,ms over all moving boxes,
\dstarAllTau$\times$ the measured interval, and \dstarMid\,ms in the
\mapCutMid--\mapCutHi\,px/s stratum. Separately, on every labeled box rather than the
velocity-evaluable subset the displacement acts on, the benchmark's own cluster resampling
holds the validation-split ordering in only \rankHoldBox\,\% of replicates, its tightest
adjacent pair separated by \rankGapMin{} points against a paired standard error of
\rankGapMinSE{}. A common displacement sweep reveals one localized
inversion between \texttt{rvt-b} and S5-S in the 25--50\,px/s stratum.
Table~\ref{tab:strata} locates the
flatness in that composition: \stShareSlow\,\% of the moving boxes travel below
\mapCutLo\,px/s, where the displacement is \stPxSlow\,px, and the sweep deepens
to \mapSpanHigh{} points between \mapCutMid{} and \mapCutHi\,px/s, where it is
\stPxHigh\,px, past the labels' own annotation dispersion of
\labelSigmaGenAnn\,px (Supplement Sec.~3).
The near invariance therefore characterizes Gen1's metric resolution, not the absence of a
temporal discrepancy.

\begin{table}[t]
\centering\footnotesize
\setlength{\tabcolsep}{1.8pt}\renewcommand{\arraystretch}{0.95}
\begin{tabular}{lrrrrr}
\toprule
speed (px/s) & share (\%) & med.\ speed & centroid (px) & $\Delta$mAP & span \\
\midrule
all moving          & 100.0        & \mapSpeedMedian & \mapPxMedian & \mapCost     & \stSpanAll \\
$<$ \mapCutLo{}     & \stShareSlow & \stSpdSlow      & \stPxSlow    & \stCostSlow  & \mapSpanSlow \\
\mapCutLo--\mapCutMid{}  & \stShareMid  & \stSpdMid  & \stPxMid     & \stCostMid   & \mapSpanMid \\
\mapCutMid--\mapCutHi{}  & \stShareHigh & \stSpdHigh & \stPxHigh    & \stCostHigh  & \mapSpanHigh \\
$\geq$ \mapCutHi{}  & \stShareTop  & \stSpdTop       & \stPxTop     & \stCostTop   & \mapSpanTop \\
\bottomrule
\end{tabular}
\caption{The sweep of Fig.~\ref{fig:predictor}b by object speed over the
\stNmoving{} ground-truth boxes with centered velocity. Other boxes are ignored.
$\Delta$mAP is the cost of shifting the ground truth to the newest-window
ablation-sensitivity centroid at $-\rvtCentroidMagMeas$\,ms. Span is the mAP
range over the $-50$ to $+30$\,ms sweep. Both values are reported in points.}
\label{tab:strata}
\end{table}


\section{Supporting audit: label-acquisition timing within one exposure}
\label{sec:flicker}

\begin{figure*}[t]
\centering
\figasset{figs/fig1_dsec_exposure.pdf}{0.80\textwidth}
\caption{Published per-frame exposure width across the \dsecSeqCount{} DSEC
training sequences whose exposure files resolve: per sequence (left, log axis,
median, interquartile range and \pctBandLo--\pctBandHi\,\% band, against the
fixed \dsecFramePeriod\,ms frame period) and pooled over \dsecFrameCount{}
frames (right). \dsecPinnedFrac\,\% of frames sit at the auto-exposure ceiling
of \dsecExpMax\,\micro s and \dsecExpGapCount{} fall between the largest
non-ceiling width of \dsecExpGapLo\,\micro s and that ceiling, separating the two exposure
regimes. The \dsecCeilingSeqs{} ceiling sequences are those measured in
Sec.~\ref{sec:flicker}.}
\label{fig:exposure}
\end{figure*}

\subsection{Exposure and labels}
\label{sec:labelclock}

The predictor-side analysis concerns the temporal support of released detector
inputs and outputs. DSEC provides a separate exposure-side analysis: published
exposure intervals bound the interval within which event timing is observed, while
within-exposure event timing contains illumination-locked structure and therefore cannot be
uniquely attributed to object timing without separating that structure.
DSEC~\cite{gehrig2021dsec} publishes
per-frame exposure metadata, and the files resolve
for
\dsecSeqCount{} train sequences. The width varies by a factor of
\dsecExpRatio{} against a frame period fixed at \dsecFramePeriod\,ms, and within
one recording by up to a factor of \dsecWithinSeqRatio{}. The auto-exposure
ceiling of \dsecExpMax\,\micro s holds \dsecPinnedFrac\,\% of frames, in
\dsecCeilingSeqs{} of the \dsecSeqCount{} sequences, and the distribution is not
continuous over that range (Fig.~\ref{fig:exposure}). The gap separates the two
exposure regimes, and ceiling exposure forms a distinct stratum.

The released DSEC-Det labels comprise \ddBoxes{} boxes on \ddTracks{} tracks
across \ddSeqs{} sequences. Every label sits on the frame clock, and the median
spacing between consecutive label times is \ddSpacingMedian\,\micro s. Over
\ddTriples{} triples of consecutive labels on a track,
\ddInterpBelowThresh\,\% of middle box centers fall within a hundredth of a
pixel of their neighbors' linear interpolant. This rules out exact local linear
interpolation as a general description of the released trajectories without
identifying the underlying annotation process. DSEC-Det ships more than one annotation
version~\cite{beyondduality2026}, the supplement names the one measured here.
Median box-center image speed is \labelSpeedMedian\,px/s, with
\labelSpeedPninetynine\,px/s at the ninety-ninth percentile.

\subsection{Mains-frequency signature and the per-pixel phase test}
\label{sec:line}

\begin{figure}[t]
\centering
\figasset{figs/fig7_ceiling_vs_day.pdf}{0.66\columnwidth}
\caption{Top: 350\,\micro s slices of the ceiling-exposure sequence
\texttt{zurich\_city\_09\_a} and of the daytime control, at the peak and trough of the
ceiling sequence's own event rate inside its densest published exposure window.
Bottom: that rate in 350\,\micro s bins, each sequence normalized to its own mean. The
ceiling rate moves by a factor of \ceilVarNight{} between its extremes, the daytime rate
by \ceilVarDay{}.}
\label{fig:flickerreal}
\end{figure}

Swiss mains is \mainsHz\,Hz, and rectified lighting can produce the tested
\flickerLineHz\,Hz modulation. The pre-specified 100-Hz line and its harmonics appear in
the event-rate spectrum of every ceiling sequence, while the daytime control lacks the
same pattern (Fig.~\ref{fig:flickerreal}). Per-pixel phase tests and the pre-specified
frequency control establish the same illumination signature.

Per pixel, the Rayleigh statistic $Z=R^{2}/N$ is $\mathrm{Exp}(1)$ under
independent uniform phases. The off-frequency and daytime controls empirically
calibrate this nominal null, while the ceiling sequence has median
$Z=\flickerZmedian{}$ at \flickerLineHz\,Hz and
\flickerFracLockedPct\,\% of live pixels reach the nominal $p<10^{-3}$ threshold under
the independent-uniform-phase reference. The illumination attribution therefore relies
primarily on the frequency-matched and daytime controls rather than on the nominal Rayleigh
reference alone.

\subsection{All six ceiling sequences, per labeled box}

The same statistic taken per labeled box, over the events inside that box during that
frame's own exposure window, extends to every ceiling sequence. Across the
\dsecCeilingSeqs{} sequences and \sixBoxes{} boxes, \sixLockedLoPct{}--\sixLockedHiPct\%
satisfy the pre-specified criterion in Supplement Table~2.


\subsection{Within-box temporal dispersion under that signature}
\label{sec:cost}

That signature limits attribution of within-exposure event timing to object timing. The
per-box 137-Hz arm is used as a same-event descriptive control rather than as an independent
spectral-resolution test. The illumination attribution is supported primarily by the long-slice
100-Hz harmonic structure and the daytime control.
Within-box dispersion is \sigmaEvtNight\,\micro s against a null of
\sigmaEvtNullNight\,\micro s, with an excess of \sigmaEvtExcessNight\,\micro s.
A frame-level permutation preserving pooled temporal distributions and box counts gives a
null of \permNull\,\micro s, and the excess remains after the clean-pixel control.
Random-position boxes show that the dispersion is not specific to annotated object regions
(Supplement Sec.~9), reinforcing an attribution-limit interpretation rather than an
object-motion interpretation.

\section{Scope and limitations}

Each result is conditional on a stated instrument: local window-ablation sensitivity,
nominal-time association under the reported regression specification, and a local
constant-velocity displacement surrogate. The SSM chunk-boundary result applies to the two
released SSM-ViT checkpoints under their released streaming evaluator. Gen1 and DSEC are
separate case studies, and no cross-dataset universality is claimed. All Gen1 quantitative claims
are validation-split claims obtained from ported released checkpoints, and sequence bootstrap
claims use the 406 sequences represented in the stored-detection evaluation.

\section{Conclusion}

In the released SSM-ViT pipeline one pooled score combines predictions with one to twenty-one
windows of recurrent history, and moving only the chunk boundary shifts S5-B by
\ssmPairedMapAllBase{} mAP over all labeled validation boxes and \ssmPairedMapVelBase{} mAP on
the velocity-evaluable subset, at fixed frames, weights, labels and evaluator. Effective
temporal support is therefore a free variable of the reported score.

The result does not single out a canonical temporal-support statistic for recurrent event
models, nor attribute the shift to one stage of the pipeline. RVT's input sensitivity and its
regression-based temporal coefficient remain distinct descriptors rather than estimates of
recurrent support, and DSEC's within-exposure timing carries illumination-locked structure that
must be separated before attribution to object motion.

Two scores at the same nominal timestamp are thus not performance at the same effective time
when the recurrent support behind them differs. Comparisons across recurrent architectures
should report or control the temporal support induced by preprocessing and recurrence---
predictor initialization, reset policy, chunk boundaries, effective history and
label-acquisition intervals---alongside the score, which makes these quantities
protocol-identifiable.
{\footnotesize
\setlength{\itemsep}{0.4pt}
\clearpage
\begin{thebibliography}{99}

\bibitem{hung2022let}
W.-C. Hung, V. Casser, H. Kretzschmar, J.-J. Hwang, and D. Anguelov.
\newblock LET-3D-AP: Longitudinal error tolerant 3D average precision for
camera-only 3D detection.
\newblock {\em arXiv:2206.07705}, 2022. Waymo Open Dataset 3D camera-only
detection challenge metric.

\bibitem{waymolabelproto}
Waymo Open Dataset.
\newblock \texttt{label.proto}, field \texttt{camera\_synced\_box} and its
documentation comment, read 2026.

\bibitem{sun2020waymo}
P. Sun et~al.
\newblock Scalability in perception for autonomous driving: Waymo Open Dataset.
\newblock In {\em IEEE/CVF Conference on Computer Vision and Pattern Recognition
(CVPR)}, 2020.

\bibitem{nunez2026annotations}
A.~J.~Miro, L.~af Klinteberg, B.~Timus, A.~Asefaw, A.~Khoche, T.~Gustafsson,
S.~S.~Mansouri, and M.~Daneshtalab.
\newblock Correcting and quantifying systematic errors in 3D box annotations for
autonomous driving.
\newblock In {\em IEEE/CVF Winter Conference on Applications of Computer Vision
(WACV)}, 2026.

\bibitem{qin2018temporal}
T. Qin and S. Shen.
\newblock Online temporal calibration for monocular visual-inertial systems.
\newblock In {\em IEEE/RSJ International Conference on Intelligent Robots and
Systems (IROS)}, 2018.

\bibitem{li2020streaming}
M. Li, Y.-X. Wang, and D. Ramanan.
\newblock Towards streaming perception.
\newblock In {\em European Conference on Computer Vision (ECCV)}, 2020.

\bibitem{mao2019delay}
H. Mao, X. Yang, and W. J. Dally.
\newblock A delay metric for video object detection: What average precision
fails to tell.
\newblock In {\em IEEE/CVF International Conference on Computer Vision (ICCV)},
2019, pages 573--582.

\bibitem{bolya2020tide}
D. Bolya, S. Foley, J. Hays, and J. Hoffman.
\newblock TIDE: A general toolbox for identifying object detection errors.
\newblock In {\em European Conference on Computer Vision (ECCV)}, 2020.

\bibitem{hoiem2012diagnosing}
D. Hoiem, Y. Chodpathumwan, and Q. Dai.
\newblock Diagnosing error in object detectors.
\newblock In {\em European Conference on Computer Vision (ECCV)}, 2012.

\bibitem{alwassel2018detad}
H. Alwassel, F. Caba Heilbron, V. Escorcia, and B. Ghanem.
\newblock Diagnosing error in temporal action detectors.
\newblock In {\em European Conference on Computer Vision (ECCV)}, 2018.

\bibitem{luiten2021hota}
J. Luiten, A. O\v{s}ep, P. Dendorfer, P. Torr, A. Geiger, L. Leal-Taix\'e, and
B. Leibe.
\newblock HOTA: A higher order metric for evaluating multi-object tracking.
\newblock {\em International Journal of Computer Vision}, 2021.

\bibitem{gehrig2021dsec}
M. Gehrig, W. Aarents, D. Gehrig, and D. Scaramuzza.
\newblock DSEC: A stereo event camera dataset for driving scenarios.
\newblock {\em IEEE Robotics and Automation Letters}, 2021.

\bibitem{gehrig2024lowlatency}
D. Gehrig and D. Scaramuzza.
\newblock Low-latency automotive vision with event cameras.
\newblock {\em Nature}, 629:1034--1040, 2024.

\bibitem{pang2021qdtrack}
J. Pang, L. Qiu, X. Li, H. Chen, Q. Li, T. Darrell, and F. Yu.
\newblock Quasi-dense similarity learning for multiple object tracking.
\newblock In {\em IEEE/CVF Conference on Computer Vision and Pattern Recognition
(CVPR)}, 2021.

\bibitem{beyondduality2026}
K. Wang, S. Liu, H. Shi, L. Shi, and H. Chen.
\newblock Beyond duality: A hybrid framework of leveraging shared and private
features for RGB-Event object detection.
\newblock In {\em IEEE/CVF Conference on Computer Vision and Pattern Recognition
(CVPR)}, 2026.

\bibitem{detournemire2020gen1}
P. de Tournemire, D. Nitti, E. Perot, D. Migliore, and A. Sironi.
\newblock A large scale event-based detection dataset for automotive.
\newblock {\em arXiv:2001.08499}, 2020.

\bibitem{perot2020megapixel}
E. Perot, P. de Tournemire, D. Nitti, J. Masci, and A. Sironi.
\newblock Learning to detect objects with a 1 megapixel event camera.
\newblock In {\em Advances in Neural Information Processing Systems (NeurIPS)},
2020.

\bibitem{cho2025ev3dod}
H. Cho, J.-Y. Kang, Y. Kim, and K.-J. Yoon.
\newblock Ev-3DOD: Pushing the temporal boundaries of 3D object detection with
event cameras.
\newblock In {\em IEEE/CVF Conference on Computer Vision and Pattern Recognition
(CVPR)}, 2025.

\bibitem{sun2023refid}
L. Sun, C. Sakaridis, J. Liang, Q. Jiang, K. Yang, P. Sun, Y. Ye, K. Wang, and
L. Van~Gool.
\newblock Event-based fusion for motion deblurring with cross-modal attention.
\newblock In {\em IEEE/CVF Conference on Computer Vision and Pattern Recognition
(CVPR)}, 2023.

\bibitem{kim2022unknownexposure}
T. Kim, J. Lee, L. Wang, and K.-J. Yoon.
\newblock Event-guided deblurring of unknown exposure time videos.
\newblock In {\em European Conference on Computer Vision (ECCV)}, 2022.

\bibitem{gehrig2023rvt}
M. Gehrig and D. Scaramuzza.
\newblock Recurrent vision transformers for object detection with event cameras.
\newblock In {\em IEEE/CVF Conference on Computer Vision and Pattern Recognition
(CVPR)}, 2023.

\bibitem{liu2024nighttime}
H. Liu, S. Peng, L. Zhu, Y. Chang, H. Zhou, and L. Yan.
\newblock Seeing Motion at Nighttime with an Event Camera.
\newblock In {\em IEEE/CVF Conference on Computer Vision and Pattern Recognition
(CVPR)}, 2024.

\bibitem{zubic2024ssm}
N. Zubi\'c, M. Gehrig, and D. Scaramuzza.
\newblock State space models for event cameras.
\newblock In {\em IEEE/CVF Conference on Computer Vision and Pattern Recognition
(CVPR)}, 2024.

\bibitem{wu2024leod}
Z. Wu, M. Gehrig, Q. Lyu, X. Liu, and I. Gilitschenski.
\newblock LEOD: Label-efficient object detection for event cameras.
\newblock In {\em IEEE/CVF Conference on Computer Vision and Pattern Recognition
(CVPR)}, pages 16933--16943, 2024.

\bibitem{sui2026astw}
Z. Sui, H. Hao, W. He, S.-H. Lee, and W. Wang.
\newblock Adaptive spatial-temporal window: unlocking the potential of event
cameras in heterogeneous velocity scenarios.
\newblock In {\em IEEE/CVF Conference on Computer Vision and Pattern Recognition
(CVPR)}, pages 946--955, 2026.

\bibitem{yang2024latency}
Y. Yang, J. Liang, B. Yu, Y. Chen, J. S. Ren, and B. Shi.
\newblock Latency correction for event-guided deblurring and frame interpolation.
\newblock In {\em IEEE/CVF Conference on Computer Vision and Pattern Recognition
(CVPR)}, pages 24977--24986, 2024.

\bibitem{chu2026stare}
J. Chu, R. Zhang, C. Yang, Z. Yu, Z. Bu, H. Liu, F. R\"ohrbein, A. Knoll,
G. Chen, and C. Jiang.
\newblock Bridging the latency gap with a continuous stream evaluation framework
in event-driven perception.
\newblock {\em Nature Communications}, 17:2441, 2026.

\bibitem{oncalibration2024}
S.~Kuzucu, K.~Oksuz, J.~Sadeghi, and P.~K.~Dokania.
\newblock On calibration of object detectors: Pitfalls, evaluation and baselines.
\newblock In {\em European Conference on Computer Vision (ECCV)}, 2024.

\bibitem{stratification2024}
R.~Fogliato, P.~Patil, M.~Monfort, and P.~Perona.
\newblock A framework for efficient model evaluation through stratification, sampling, and
estimation.
\newblock In {\em European Conference on Computer Vision (ECCV)}, 2024.

\bibitem{northcutt2021labelerrors}
C. G. Northcutt, A. Athalye, and J. Mueller.
\newblock Pervasive label errors in test sets destabilize machine learning
benchmarks.
\newblock In {\em NeurIPS Datasets and Benchmarks}, 2021.

\bibitem{ercan2023evreal}
B. Ercan, O. Eker, A. Erdem, and E. Erdem.
\newblock EVREAL: Towards a comprehensive benchmark and analysis suite for
event-based video reconstruction.
\newblock In {\em IEEE/CVF Conference on Computer Vision and Pattern Recognition
Workshops (CVPRW)}, 2023.

\bibitem{yang2022streamyolo}
J. Yang, S. Liu, Z. Li, X. Li, and J. Sun.
\newblock Real-time object detection for streaming perception.
\newblock In {\em IEEE/CVF Conference on Computer Vision and Pattern Recognition
(CVPR)}, 2022.

\bibitem{softed2024}
R. Salles et~al.
SoftED: Metrics for soft evaluation of time series event detection.
\newblock {\em Computers \& Industrial Engineering}, 2024.

\bibitem{faod2024}
Haitian Zhang et~al.
Frequency-adaptive low-latency object detection using events and frames.
\newblock {\em arXiv:2412.04149}, 2024. Unrefereed preprint.

\end{thebibliography}
}

\end{document}
